\frfilename{mt282.tex}
\versiondate{24.9.09}
\copyrightdate{1996}

\def\chaptername{Analisis Fourier}
\def\sectionname{Deret Fourier}
\def\Var{\mathop{\text{Var}}}

\newsection{282}

Dari sekian banyak teori deret Fourier, saya mengekstrak beberapa hasil
yang setidaknya dapat memberikan dasar untuk studi lebih lanjut.   Saya memberikan definisi
dari jumlah Fourier dan Fej\'er (282A), dengan lima hasil
yang paling penting mengenai konvergensinya (282G, 282H, 282J, 282L,
282O).   Dalam perjalanan saya menyertakan lema Riemann-Lebesgue (282E).   Saya mengakhiri
dengan menyebutkan konvolusi (282Q).

\vleader{48pt}{282A}{Definisi} Misalkan $f$ adalah fungsi bernilai kompleks
yang dapat diintegralkan dan terdefinisi hampir di mana-mana pada $\ocint{-\pi,\pi}$.

\header{282Aa}{\bf (a)} {\bf Koefisien Fourier} dari $f$ adalah bilangan kompleks

$$c_k=\Bover1{2\pi}\int_{-\pi}^{\pi}f(x)e^{-ikx}dx$$

\noindent untuk $k\in\Bbb Z$.

\header{282Ab}{\bf (b)} {\bf Jumlah Fourier} dari $f$ adalah fungsi-fungsi

$$s_n(x)=\sum_{k=-n}^nc_ke^{ikx}$$

\noindent untuk $x\in\ocint{-\pi,\pi}$, $n\in\Bbb N$.

\header{282Ac}{\bf (c)} {\bf Deret Fourier
} dari $f$ adalah deret $\sum_{k=-\infty}^{\infty}c_ke^{ikx}$,
atau\cmmnt{ (karena
biasanya kita meninjau jumlah parsial simetris $s_n$)} deret
$c_0+\sum_{k=1}^{\infty}(c_ke^{ikx}+c_{-k}e^{-ikx})$.

\header{282Ad}{\bf (d)} {\bf Jumlah Fej\'er} dari $f$ adalah fungsi-fungsi

$$\sigma_m=\Bover1{m+1}\sum_{n=0}^ms_n$$

\noindent untuk $m\in\Bbb N$.

\header{282Ae}{\bf (e)}\cmmnt{ Akan berguna jika kita memiliki satu istilah
lagi.}   Jika $f$ adalah
sembarang fungsi dengan $\dom f\subseteq\ocint{-\pi,\pi}$, {\bf perluasan
periodik}-nya adalah fungsi $\tilde f$, dengan domain
$\bigcup_{k\in\Bbb Z}(\dom f+2k\pi)$, sehingga $\tilde f(x)=f(x-2k\pi)$
kapan pun $k\in\Bbb Z$ dan $x\in\dom f+2k\pi$.

\cmmnt{
\leader{282B}{Keterangan} Di sini saya telah membuat dua pilihan yang kurang lebih sewenang-wenang.

\header{282Ba}{\bf (a)} Saya telah memilih untuk mengekspresikan deret Fourier dalam bentuk `kompleks'
daripada
dalam bentuk `real'.   Dari sudut pandang teori ukuran murni
(dan memang dari sudut pandang asal-usul bidang ini
pada abad kesembilan belas), ada keuntungan dalam keanggunan bila perhatian diarahkan pada fungsi bernilai real
$f$ dan koefisien real

$$a_k=\Bover1{\pi}\int_{-\pi}^{\pi}f(x)\cos kx\,dx\text{ untuk }k\in\Bbb
N,$$

$$b_k=\Bover1{\pi}\int_{-\pi}^{\pi}f(x)\sin kx\,dx\text{ untuk }k\ge
1.$$

\noindent Jika kita memilih bentuk ini, kita mempunyai

\Centerline{$c_0=\Bover12a_0$,}

\noindent dan untuk $k\ge 1$ kita mempunyai

\Centerline{$c_k=\Bover12(a_k-ib_k)$,\quad $c_{-k}=\Bover12(a_k+ib_k)$,
\quad
$a_k=c_k+c_{-k}$,\quad $b_k=i(c_k-c_{-k})$,}

\noindent sehingga jumlah Fourier menjadi

$$s_n(x)=\Bover12a_0+\sum_{k=1}^na_k\cos kx+b_k\sin kx.$$

\noindent Keuntungannya adalah bahwa fungsi bernilai real $f$ bersesuaian dengan koefisien real
$a_k$, $b_k$, sehingga jelas bahwa jika $f$ bernilai real
maka jumlah Fourier dan Fej\'er juga bernilai real.   Kerugiannya adalah bahwa kita harus
menggunakan berbagai persamaan trigonometri yang lebih rumit
daripada sifat-sifat fungsi eksponensial kompleks
yang dicerminkannya, dan kita semakin jauh dari generalisasi
alami ke grup abelian kompak lokal.   Jadi baik insinyur kelistrikan
maupun analis harmonik cenderung lebih menyukai koefisien $c_k$.

\header{282Bb}{\bf (b)} Saya memilih fungsi $f$ yang terdefinisi pada
interval $\ocint{-\pi,\pi}$ daripada pada lingkaran
$S^1=\{z:z\in\Bbb C,\,|z|=1\}$.   Akan ada keuntungan dalam keanggunan bahasa
dalam menggunakan $S^1$, meskipun saya tidak ingat sering melihat rumus


\Centerline{$c_k=\int z^kf(z)dz$}

\noindent yang merupakan terjemahan alami dari
$c_k=\bover1{2\pi}\int e^{ikx}f(x)dx$ di bawah substitusi $x=\arg z$,
$dx=2\pi\nu(dz)$.   Namun, penerapan teori ini cenderung berurusan
dengan fungsi periodik pada garis real, jadi saya bekerja dengan
$\ocint{-\pi,\pi}$, dan menerima kenyataan bahwa operasi grupnya
$+_{2\pi}$, menulis
$x+_{2\pi}y$ untuk mana pun di antara $x+y$, $x+y+2\pi$, $x+y-2\pi$ yang termasuk dalam
$\ocint{-\pi,\pi}$, kurang familiar dibandingkan perkalian pada
$S^1$.

\header{282Bc}{\bf (c)} Pernyataan di (b) seharusnya mengingatkan Anda pada
\S255.

\header{282Bd}{\bf (d)} Perhatikan bahwa jika $f\eae g$ maka $f$ dan $g$ memiliki
koefisien Fourier, jumlah Fourier, dan jumlah Fej\'er yang sama.
Ini berarti bahwa kita dapat, jika mau, menganggap $c_k$, $s_n$ dan
$\sigma_m$ dikaitkan dengan anggota $L_{\Bbb C}^1$, ruang kelas kesetaraan
dari fungsi yang dapat diintegralkan (\S242), bukan sebagai
yang dikaitkan dengan fungsi tertentu $f$.   Karena $s_n$ dan
$\sigma_m$ muncul sebagai fungsi sebenarnya, dan karena banyak pertanyaan
yang kami minati mengacu pada nilainya pada titik tertentu, maka
lebih wajar untuk menyatakan teori dalam istilah fungsi yang dapat diintegralkan $f$
daripada dalam istilah anggota $L_{\Bbb C}^1$.
}%end of comment

\leader{282C}{Masalah (a)} Dalam kondisi apa, dan dalam pengertian apa, jumlah Fourier dan Fej\'er $s_n$ dan $\sigma_m$ dari fungsi
$f$ konvergen ke $f$?

\header{282Cb}{\bf (b)} Bagaimana properti barisan
berujung ganda $\langle c_k\rangle_{k\in\Bbb Z}$ mencerminkan properti
$f$, dan sebaliknya?

\cmmnt{\medskip

\noindent{\bf Catatan} Teori deret Fourier telah menjadi salah satu
topik utama analisis selama hampir dua ratus tahun, dan tak terhitung banyaknya
masalah lebih lanjut telah memberikan kontribusi besar terhadap pemahaman kita.   (Untuk
contoh: dapatkah kita mencirikan barisan $\langle
c_k\rangle_{k\in\Bbb Z}$ yang merupakan koefisien Fourier dari suatu
fungsi yang dapat diintegralkan?) Namun dalam garis besar ini saya akan berkonsentrasi pada pertanyaan (a)
di atas, dengan hasil satu setengah (282K, 282Rb) mengatasi (b), yang
akan memberi kita lebih dari cukup bahan untuk dikerjakan.

Meskipun kebanyakan orang akan merasa bahwa jumlah Fourier lebih dekat dengan
apa yang sebenarnya ingin kita ketahui, ternyata jumlah Fej\'er lebih mudah dianalisis, dan ada keuntungan jika menanganinya terlebih dahulu.
Jadi, meskipun Anda mungkin ingin melihat pernyataan 282J, 282L dan
282O untuk mendapatkan gambaran tentang tujuan kita, paruh pertama bagian ini
sebagian besar akan membahas jumlah Fej\'er.   Perhatikan bahwa dalam setiap kasus ketika kita
tahu bahwa jumlah Fourier konvergen (hal ini cukup umum; lihat, misalnya, contoh di 282Xh dan 282Xo), jika kita tahu bahwa jumlah
Fej\'er konvergen ke $f$, kita dapat menyimpulkan dari 273Ca bahwa jumlah Fourier juga demikian.

Langkah pertama adalah lema dasar yang menunjukkan bahwa jumlah Fourier dan
Fej\'er dari suatu fungsi $f$ dapat dianggap sebagai konvolusi $f$
dengan kernel yang dapat dinyatakan dengan fungsi-fungsi yang dikenal.
}%end of comment

\leader{282D}{Lema} Misalkan $f$ adalah fungsi bernilai kompleks yang
dapat diintegralkan pada $\ocint{-\pi,\pi}$, dan

\Centerline{$c_k=\bover1{2\pi}\int_{-\pi}^{\pi}f(x)e^{-ikx}dx$,
\quad$s_n(x)=\sum_{k=-n}^nc_ke^{ikx}$,
\quad$\sigma_m(x)=\bover1{m+1}\sum_{n=0}^ms_n(x)$}

\noindent masing-masing adalah koefisien Fourier, jumlah Fourier, dan jumlah Fej\'er-nya.
Tulis $\tilde f$ untuk perpanjangan periodik $f$\cmmnt{ (282Ae)}.   Untuk
$m\in\Bbb N$, tulis

$$\psi_m(t)=\Bover{1-\cos(m+1)t}{2\pi(m+1)(1-\cos t)}$$

\noindent untuk $0<|t|\le\pi$.   \cmmnt{(Jika mau, Anda dapat mengatur
$\psi_m(0)=\bover{m+1}{2\pi}$ agar $\psi_m$ kontinu di
$[-\pi,\pi]$.)}

(a) Untuk setiap $n\in\Bbb N$, $x\in\ocint{-\pi,\pi}$,

$$\eqalign{s_n(x)
&=\Bover1{2\pi}\int_{-\pi}^{\pi}f(t)
\Bover{\sin(n+{1\over 2}\pushbottom{3.5pt})(x-t)}
  {\sin{1\over 2}(x-t)}dt\cr
&=\Bover1{2\pi}\int_{-\pi}^{\pi}\tilde f(x+t)
\Bover{\sin(n+{1\over 2}\pushbottom{3.5pt})t}{\sin{1\over 2}t}dt\cr
&=\Bover1{2\pi}\int_{-\pi}^{\pi}f(x-_{2\pi}t)
\Bover{\sin(n+{1\over 2}\pushbottom{3.5pt})t}{\sin{1\over 2}t}dt,\cr}$$

\noindent menulis $x-_{2\pi}t$ untuk $x-t$, $x-t-2\pi$,
$x-t+2\pi$ mana pun yang termasuk dalam $\ocint{-\pi,\pi}$.

(b) Untuk setiap $m\in\Bbb N$, $x\in\ocint{-\pi,\pi}$,

$$\eqalign{\sigma_m(x)
&=\int_{-\pi}^{\pi}\tilde f(x+t)\psi_m(t)dt\cr
&=\int_{0}^{\pi}(\tilde f(x+t)+\tilde f(x-t))\psi_m(t)dt\cr
&=\int_{-\pi}^{\pi}f(x-_{2\pi}t)\psi_m(t)dt.\cr}$$

(c) Untuk setiap $n\in\Bbb N$,

$$\Bover1{2\pi}\int_{-\pi}^{0}
  \Bover{\sin(n+{1\over 2}\pushbottom{3.5pt})t}{\sin{1\over 2}t}dt
=\Bover1{2\pi}\int_{0}^{\pi}
\Bover{\sin(n+{1\over 2}\pushbottom{3.5pt})t}{\sin{1\over 2}t}dt
=\Bover12,
\quad\Bover1{2\pi}\int_{-\pi}^{\pi}
\Bover{\sin(n+{1\over 2}\pushbottom{3.5pt})t}{\sin{1\over 2}t}dt
=1.$$

(d) Untuk setiap $m\in\Bbb N$,

\quad (i) $0\le\psi_m(t)\le\Bover{m+1}{2\pi}$ untuk setiap $t$;

\quad (ii) untuk $\delta>0$, $\lim_{m\to\infty}\psi_m(t)=0$ secara seragam
di $\{t:\delta\le|t|\le\pi\}$;

\quad (iii) $\int_{-\pi}^0\psi_m=\int_0^{\pi}\psi_m=\Bover12$,
\quad $\int_{-\pi}^{\pi}\psi_m=1$.


\proof{ Sebenarnya semua ini hanyalah menjumlahkan deret geometri.

\medskip

{\bf (a)} Untuk (a), kita punya

$$\eqalign{\sum_{k=-n}^ne^{-ikt}
&=\bover{e^{int}-e^{-i(n+1)t}}{1-e^{-it}}\cr
&=\bover{e^{i(n+{1\over 2})t}-e^{-i(n+{1\over 2})t}}
{e^{{1\over 2}it}-e^{-{1\over 2}it}}
=\bover{\sin(n+{1\over 2}\pushbottom{3.5pt})t}{\sin{1\over 2}t}.\cr}$$

\noindent Jadi

$$\eqalign{s_n(x)
&=\sum_{k=-n}^nc_ke^{ikx}
=\Bover1{2\pi}\int_{-\pi}^{\pi}f(t)
  \bigl(\sum_{k=-n}^ne^{ik(x-t)}\bigr)dt\cr
&=\Bover1{2\pi}\int_{-\pi}^{\pi}f(t)
  \Bover{\sin(n+{1\over 2}\pushbottom{3.5pt})(x-t)}
     {\sin{1\over 2}(x-t)}dt
=\Bover1{2\pi}\int_{-\pi}^{\pi}\tilde f(t)
  \Bover{\sin(n+{1\over 2}\pushbottom{3.5pt})(x-t)}
     {\sin{1\over 2}(x-t)}dt\cr
&=\Bover1{2\pi}\int_{-\pi-x}^{\pi-x}\tilde f(x+t)
  \Bover{\sin(n+{1\over 2}\pushbottom{3.5pt})t}{\sin{1\over 2}t}dt
=\Bover1{2\pi}\int_{-\pi}^{\pi}\tilde f(x+t)
  \Bover{\sin(n+{1\over 2}\pushbottom{3.5pt})t}{\sin{1\over 2}t}dt\cr}$$

\noindent karena $\tilde f$ dan $t\mapsto
\bover{\sin(n+{1\over 2})t}{\sin{1\over 2}t}$ periodik dengan periode
$2\pi$, sehingga integral dari $-\pi-x$ ke $-\pi$ harus sama dengan
integral dari $\pi-x$ ke $\pi$.

Untuk ekspresi dalam $f(x-_{2\pi}t)$, kita punya

$$\eqalignno{s_n(x)
&=\Bover1{2\pi}\int_{-\pi}^{\pi}\tilde f(x+t)
\Bover{\sin(n+{1\over 2}\pushbottom{3.5pt})t}{\sin{1\over 2}t}dt
=\Bover1{2\pi}\int_{-\pi}^{\pi}\tilde f(x-t)
\Bover{\sin(n+{1\over 2}\pushbottom{3.5pt})(-t)}
   {\sin{1\over 2}(-t)}dt\cr
\noalign{\noindent (dengan mengganti $-t$ untuk $t$)}
&=\Bover1{2\pi}\int_{-\pi}^{\pi}f(x-_{2\pi}t)
\Bover{\sin(n+{1\over 2}\pushbottom{3.5pt})t}{\sin{1\over 2}t}dt\cr}$$

\noindent karena (untuk $x$, $t\in\ocint{-\pi,\pi}$)
$f(x-_{2\pi}t)=\tilde f(x-t)$ setiap kali
didefinisikan, dan $\sin$ merupakan fungsi ganjil.

\medskip

{\bf (b)} Dengan cara yang sama, kita punya

$$\eqalign{\sum_{n=0}^m\sin(n+\Bover12)t
&=\Imag\bigl(\sum_{n=0}^me^{i(n+{1\over 2})t}\bigr)
=\Imag\bigl(e^{{1\over 2}it}\sum_{n=0}^me^{int}\bigr)\cr
&=\Imag\bigl(e^{{1\over 2}it}\bover{1-e^{i(m+1)t}}{1-e^{it}}\bigr)
=\Imag\bigl(\bover{1-e^{i(m+1)t}}{e^{-{1\over 2}it}
   -e^{{1\over2}it}}\bigr)\cr
&=\Imag\bigl(\bover{1-e^{i(m+1)t}}{-2i\sin{1\over 2}t}\bigr)
=\Imag\bigl(\bover{i(1-e^{i(m+1)t})}{2\sin{1\over 2}t}\bigr)\cr
&=\bover{1-\cos(m+1)t}{2\sin{1\over 2}t}.\cr}$$

\noindent Jadi

\Centerline{$\sum_{n=0}^m
  \Bover{\sin(n+{1\over 2}\pushbottom{3.5pt})t}{\sin{1\over 2}t}
=\Bover{1-\cos(m+1)t}{2\sin^2{1\over 2}t}
=\Bover{1-\cos(m+1)t}{1-\cos t}
=2\pi(m+1)\psi_m(t)$.}

\noindent Oleh karena itu,

$$\eqalignno{\sigma_m(x)
&={\Bover1{m+1}}\sum_{n=0}^ms_n(x)\cr
&={\Bover1{m+1}}\sum_{n=0}^m{\Bover1{2\pi}}\int_{-\pi}^{\pi}
   \tilde f(x+t){\Bover{\sin(n+{1\over
2}\pushbottom{3.5pt})t}{\sin{1\over 2}t}}dt\cr
&={\Bover1{2\pi}}\int_{-\pi}^{\pi}
   \tilde f(x+t)\bigl({\Bover1{m+1}}\sum_{n=0}^m
   {\Bover{\sin(n+{1\over 2}\pushbottom{3.5pt})t}{\sin{1\over 2}t}}
     \bigr)dt\cr
&=\int_{-\pi}^{\pi}\tilde f(x+t)\psi_m(t)dt
=\int_{-\pi}^{\pi}f(x-_{2\pi}t)\psi_m(t)dt\cr}$$

\noindent seperti pada (a), karena $\cos$ dan $\psi_m$ merupakan fungsi genap.
Untuk alasan yang sama,

$$\int_{0}^{\pi}\tilde f(x-t)\psi_m(t)dt
=\int_{-\pi}^{0}\tilde f(x+t)\psi_m(t)dt,$$

\noindent Jadi

$$\sigma_m(x)=\int_0^{\pi}(\tilde f(x+t)+\tilde f(x-t))\psi_m(t)dt.$$


\medskip

{\bf (c)} Kita hanya perlu melihat dari mana rumus
$\bover{\sin(n+{1\over 2})t}{\sin{1\over 2}t}$ berasal untuk melihat bahwa

$$\eqalign{\Bover1{2\pi}\int_I
  \Bover{\sin(n+{1\over 2}\pushbottom{3.5pt})t}{\sin{1\over 2}t}dt
&=\Bover1{2\pi}\int_I\sum_{k=-n}^ne^{ikt}dt\cr
&=\Bover1{2\pi}\int_I(1+2\sum_{k=1}^n\cos kt)dt
=\Bover12\cr}$$

\noindent untuk $I=[-\pi,0]$ dan $I=[0,\pi]$, karena
$\int_I\cos kt\,dt=0$ untuk setiap $k\ne 0$.

\medskip

{\bf (d)(i)} $\psi_m(t)\ge 0$ untuk setiap $t$ karena
$1-\cos(m+1)t$, $1-\cos t$ selalu lebih besar atau sama dengan $0$.
Untuk batas atas, kita punya, dengan menggunakan konstruksi pada (a) dan (b),

$$\bigl|\Bover{\sin(n+{1\over 2}\pushbottom{3.5pt})t}
   {\sin{1\over 2}t}\bigr|
=\bigl|\sum_{k=-n}^ne^{ikt}\bigr|
\le 2n+1$$

\noindent untuk setiap $n$, jadi

$$\eqalign{\psi_m(t)
&=\Bover1{2\pi(m+1)}\sum_{n=0}^m\Bover{\sin(n+{1\over
2}\pushbottom{3.5pt})t}{\sin{1\over 2}t}\cr
&\le\Bover1{2\pi(m+1)}\sum_{n=0}^m\,2n+1
=\Bover{m+1}{2\pi}.\cr}$$

\medskip

\quad{\bf (ii)} Jika $\delta\le|t|\le\pi$,

\Centerline{$\psi_m(t)\le\Bover1{\pi(m+1)(1-\cos t)}
\le\Bover1{\pi(m+1)(1-\cos\delta)}\to 0$}

\noindent sebagai $m\to\infty$.

\medskip

\quad{\bf (iii)} juga mengikuti konstruksi pada (b), karena

$$\int_I\psi_m
={\Bover1{2\pi(m+1)}}\sum_{n=0}^m
  \int_I\Bover{\sin(n+{1\over 2}\pushbottom{3.5pt})t}{\sin{1\over 2}t}dt
={\Bover1{m+1}}\sum_{n=0}^m\Bover12
=\Bover12$$

\noindent untuk $I=[-\pi,0]$ dan $I=[0,\pi]$, menggunakan (c).
}%end of proof of 282D

\cmmnt{\medskip

\noindent{\bf Keterangan} Untuk pembahasan substitusi integral, jika
Anda merasa perlu membenarkan manipulasi pada bagian (a) pembuktian,
lihat 263J.

Fungsi-fungsi

\Centerline{$t\mapsto\Bover{\sin(n+{1\over
2}\pushbottom{3.5pt})t}{\sin{1\over 2}t}$,
\quad $t\mapsto \Bover{1-\cos(m+1)t}{(m+1)(1-\cos t)}$}

\noindent disebut masing-masing kernel {\bf Dirichlet} dan kernel
{\bf Fej\'er}.

Rumus $f(x-_{2\pi}t)$ saya berikan pada (a) dan (b) agar
memberikan rujukan alami ke pembahasan dalam 255O.
}

\leader{282E}{}\cmmnt{ Langkah selanjutnya adalah lema penting, dengan nama khusus
yang sesuai (Anda akan senang mengetahuinya) mencerminkan pentingnya,
bukan kesulitannya.

\medskip

\noindent}{\bf Lema Riemann-Lebesgue} Misalkan $f$ adalah fungsi bernilai kompleks
yang dapat diintegralkan pada $\Bbb R$.   Maka

\Centerline{$\lim_{y\to\infty}\int f(x)e^{-iyx}dx
=\lim_{y\to-\infty}\int f(x)e^{-iyx}dx
=0$.}

\proof{{\bf (a)} Pertimbangkan dulu kasus di mana $f=\chi\ooint{a,b}$,
di mana $a<b$.   Maka

\Centerline{$|\int f(x)e^{-iyx}dx|
=|\int_a^{b}e^{-iyx}dx|
=|\Bover1{-iy}(e^{-iyb}-e^{-iya})|
\le\Bover2{|y|}$}

\noindent jika $y\ne 0$.   Jadi dalam hal ini hasilnya sudah jelas.

\medskip

{\bf (b)} Segera dapat disimpulkan bahwa hasilnya benar jika $f$ adalah fungsi tangga
dengan dukungan terbatas, yaitu, jika ada $a_0\le
a_1\ldots\le a_n$ sehingga
$f$ konstan pada setiap interval
$\ooint{a_{j-1},a_j}$ dan nol di luar $[a_0,a_n]$.

\medskip

{\bf (c)} Sekarang, untuk $f$ yang dapat diintegralkan dan $\epsilon>0$, terdapat
fungsi tangga $g$ sehingga $\int|f-g|\le\epsilon$
(242Oa).   Jadi

\Centerline{$|\int f(x)e^{-iyx}dx-\int g(x)e^{-iyx}dx|
\le\int|f(x)-g(x)|dx\le\epsilon$}

\noindent untuk setiap $y$, dan

\Centerline{$\limsup_{y\to\infty}|\int f(x)e^{-iyx}dx|\le\epsilon$,}

\Centerline{$\limsup_{y\to-\infty}|\int f(x)e^{-iyx}dx|\le\epsilon$.}

\noindent Karena $\epsilon$ bersifat arbitrer, kami mendapatkan hasilnya.
}%end of proof of 282E

\leader{282F}{Akibat} (a) Misalkan $f$ adalah fungsi
bernilai kompleks yang dapat diintegralkan pada $\ocint{-\pi,\pi}$, dan
$\langle c_k\rangle_{k\in\Bbb Z}$
merupakan barisan koefisien Fourier-nya. Maka
$\lim_{k\to\infty}c_k=\lim_{k\to-\infty}c_k=0$.

(b) Misalkan $f$ adalah fungsi bernilai kompleks yang dapat diintegralkan pada
$\Bbb R$.   Maka
\penalty -100
$\lim_{y\to\infty}\int f(x)\sin yx\,dx=0$.

\wheader{282F}{6}{2}{2}{36pt}

\proof{{\bf (a)} Kita hanya perlu mengidentifikasi

$$c_k=\Bover1{2\pi}\int_{-\pi}^{\pi}f(x)e^{-ikx}dx$$

\noindent dengan $\int g(x)e^{-ikx}dx$, dimana $g(x)=f(x)/2\pi$ untuk
$x\in\dom f$ dan $0$ untuk $|x|>\pi$.

\medskip

{\bf (b)} Ini hanya karena

\Centerline{$\int f(x)\sin yx\,dx
=\Bover1{2i}(\int f(x)e^{iyx}dx-\int f(x)e^{-iyx}dx)$.}
}%end of proof of 282F

\leader{282G}{}\cmmnt{ Sekarang kita siap untuk teorema tentang
konvergensi jumlah Fej\'er.
Saya mulai dengan latihan yang mudah, hampir seperti latihan pemanasan.

\medskip

\noindent}{\bf Teorema} Misalkan $f:\ocint{-\pi,\pi}\to\Bbb C$ adalah fungsi kontinu
sehingga $\lim_{t\downarrow-\pi}f(t)=f(\pi)$.
Maka urutannya $\sequence{m}{\sigma_m}$ dari jumlah
Fej\'er konvergen secara seragam ke $f$ di $\ocint{-\pi,\pi}$.

\proof{   Ketentuan pada $f$ hanya menyatakan bahwa
ekstensi periodiknya $\tilde f$ ditentukan dan kontinu di mana saja
di $\Bbb R$.   Akibatnya ia terbatas dan kontinu seragam
secara seragam pada setiap interval terbatas, khususnya pada interval $[-2\pi,2\pi]$.
Tetapkan
$K=\sup_{|t|\le 2\pi}|\tilde f(t)|=\sup_{t\in\ocint{-\pi,\pi}}|f(t)|$.
Menulis

\Centerline{$\psi_m(t)=\Bover{1-\cos(m+1)t}{2\pi(m+1)(1-\cos t)}$}

\noindent untuk $m\in\Bbb N$,  $0<|t|\le\pi$, seperti pada 282D.

Untuk setiap $\epsilon>0$ kita dapat menemukan $\delta\in\ocint{0,\pi}$ sehingga
$|\tilde f(x+t)-\tilde f(x)|\le\epsilon$ kapan pun $x\in[-\pi,\pi]$ dan
$|t|\le\delta$.   Selanjutnya, kita dapat menemukan $m_0\in\Bbb N$ sehingga
$M_m\le\bover{\epsilon}{4\pi K}$ untuk setiap $m\ge m_0$, di mana
$M_m=\sup_{\delta\le|t|\le\pi}\psi_m(t)$ (282D(d-ii)).   Sekarang anggaplah
$m\ge m_0$ dan $x\in\ocint{-\pi,\pi}$.   Tetapkan
$g(t)=\tilde f(x+t)-f(x)$ untuk $|t|\le\pi$.   Maka $|g(t)|\le 2K$ untuk semua
$t\in[-\pi,\pi]$ dan $|g(t)|\le\epsilon$ jika $|t|\le\delta$, jadi

$$\eqalign{\bigl|\int_{-\pi}^{\pi}g\times\psi_m\bigr|
&\le\int_{-\pi}^{-\delta}|g|\times\psi_m
  +\int_{-\delta}^{\delta}|g|\times\psi_m
  +\int_{\delta}^{\pi}|g|\times\psi_m\cr
&\le 2M_mK(\pi-\delta)
  +\epsilon\int_{-\delta}^{\delta}\psi_m
  +2M_mK(\pi-\delta)\cr
&\le 4\pi M_mK+\epsilon
\le 2\epsilon.\cr}$$

\noindent Akibatnya, menggunakan 282Db dan 282D(d-iii),

$$|\sigma_m(x)-f(x)|
=|\int_{-\pi}^{\pi}(\tilde f(x+t)-f(x))\psi_m(t)dt|
\le 2\epsilon$$

\noindent untuk setiap $m\ge m_0$;  dan ini berlaku untuk setiap
$x\in\ocint{-\pi,\pi}$.   Karena $\epsilon$ bersifat arbitrer,
$\sequence{m}{\sigma_m}$ konvergen ke $f$ secara seragam di
$\ocint{-\pi,\pi}$.
}%end of proof of 282G

\leader{282H}{}\cmmnt{ Sekarang saya sampai pada teorema yang menjelaskan perilaku
dari jumlah fungsi umum
Fej\'er $f$.   Hipotesis teorema
mungkin perlu sedikit dicerna;  Anda bisa mendapatkan gambaran tentang cakupan
yang dimaksudkan dengan melihat Akibat 282I.

\medskip

\noindent}{\bf Teorema} Misalkan $f$ adalah fungsi bernilai kompleks yang
dapat diintegralkan pada $\ocint{-\pi,\pi}$, dan $\sequence{m}{\sigma_m}$
barisan jumlah Fej\'er-nya. Misalkan $x\in\ocint{-\pi,\pi}$ dan
$c\in\Bbb C$ sedemikian sehingga

$$\lim_{\delta\downarrow
0}\Bover1{\delta}\int_0^{\delta}|\tilde f(x+t)+\tilde f(x-t)-2c|dt=0,$$

\noindent menulis $\tilde f$ untuk perpanjangan berkala $f$\cmmnt{, seperti
biasa};  lalu $\lim_{m\to\infty}\sigma_m(x)=c$.

\proof{ Tetapkan $\phi(t)=|\tilde f(x+t)+\tilde f(x-t)-2c|$ ketika
ini didefinisikan, yang hampir ada di mana-mana, dan
$\Phi(t)=\int_0^t\phi$, yang didefinisikan untuk setiap $t\ge 0$,
karena $\tilde f$ dapat diintegralkan pada $\ocint{-\pi,\pi}$ dan oleh karena itu
pada setiap interval yang dibatasi.

Seperti di 282D, tetapkan

\Centerline{$\psi_m(t)=\Bover{1-\cos(m+1)t}{2\pi(m+1)(1-\cos t)}$}

\noindent untuk $m\in\Bbb N$, $0<|t|\le\pi$.    Kita punya

$$|\sigma_m(x)-c|
=|\int_0^{\pi}(\tilde f(x+t)+\tilde f(x-t)-2c)\psi_m(t)dt|
\le\int_0^{\pi}\phi(t)\psi_m$$

\noindent oleh (b) dan (d) dari 282D.

Misalkan $\epsilon>0$.   Berdasarkan hipotesis, $\lim_{t\downarrow 0}\Phi(t)/t=0$;
misalkan $\delta\in\ocint{0,\pi}$ menjadi seperti $\Phi(t)\le\epsilon t$ 
setiap $t\in[0,\delta]$.   Ambil sembarang $m\ge \pi/\delta$.   Saya membagi
integral $\int_0^{\pi}\phi\times\psi_m$ menjadi tiga bagian.

\medskip

{\bf (i)} Untuk integral dari $0$ ke $1/m$, kita punya

$$\int_0^{1/m}\phi\times\psi_m
\le\int_0^{1/m}\Bover{m+1}{2\pi}\phi
=\Bover{m+1}{2\pi}\Phi(\Bover1m)
\le\Bover{\epsilon(m+1)}{2\pi m}
\le\epsilon,$$

\noindent karena $\psi_m(t)\le\bover{m+1}{2\pi}$ untuk setiap $t$
(282D(d-i)).

\medskip

{\bf (ii)} Untuk integral dari $1/m$ ke $\delta$, kita punya

$$\eqalignno{\int_{1/m}^{\delta}\phi\times\psi_m
&\le\bover1{2\pi(m+1)}\int_{1/m}^{\delta}\phi(t)\bover{1}{1-\cos t}dt
\le\bover{\pi}{4(m+1)}\int_{1/m}^{\delta}\bover{\phi(t)}{t^2}dt\cr
\noalign{\noindent (karena $1-\cos t\ge\Bover{2t^2}{\pi^2}$ untuk
$|t|\le\pi$)}
&=\bover{\pi}{4(m+1)}\bigl(\bover{\Phi(\delta)}{\delta^2}
-\bover{\Phi({1\over m})}{({1\over m})^2}
+\int_{1/m}^{\delta}\bover{2\Phi(t)}{t^3}dt\bigr)\cr
\noalign{\noindent (dengan integrasi parsial; lihat 225F)}
&\le\bover{\pi}{4(m+1)}\bigl(\bover{\epsilon}{\delta}
+\int_{1/m}^{\delta}\bover{2\epsilon}{t^2}dt\bigr)\cr
\noalign{\noindent (karena $\Phi(t)\le\epsilon t$ untuk $0\le
t\le\delta$)}
&\le\bover{\pi}{4(m+1)}\bigl(\bover{\epsilon}{\delta}
   +2\epsilon m\bigr)
\le\bover{\pi\epsilon}{4(m+1)\delta}+\bover{\pi\epsilon}2
\le\bover{\epsilon}4+\bover{\pi\epsilon}2
\le 2\epsilon.\cr}$$

\medskip

{\bf (iii)} Untuk integral dari $\delta$ ke $\pi$, kita punya

$$\int_{\delta}^{\pi}\phi\times\psi_m
\le\int_{\delta}^{\pi}\Bover1{\pi(m+1)(1-\cos\delta)}\phi
\to 0\text{ ketika }m\to\infty$$

\noindent karena $\phi$ dapat diintegralkan pada $[-\pi,\pi]$.   Oleh karena itu
ada $m_0\in\Bbb N$ sedemikian sehingga

$$\int_{\delta}^{\pi}\phi\times\psi_m\le\epsilon$$

\noindent untuk setiap $m\ge m_0$.

\medskip

Jika digabungkan, kita melihat bahwa

$$\int_0^{\pi}\phi\times\psi_m
\le\epsilon+2\epsilon
+\epsilon= 4\epsilon$$

\noindent untuk setiap $m\ge \max(m_0,\bover{\pi}{\delta})$.   Karena
$\epsilon$ bersifat arbitrer,
$\lim_{m\to\infty}\sigma_m(x)=c$, seperti yang diklaim.
}%end of proof of 282H

\leader{282I}{Akibat} Misalkan $f$ adalah fungsi bernilai kompleks yang dapat
diintegralkan pada $\ocint{-\pi,\pi}$, dan $\sequence{m}{\sigma_m}$ barisan
jumlah Fej\'er-nya.

(a) $f(x)=\lim_{m\to\infty}\sigma_m(x)$ untuk hampir setiap
$x\in\ocint{-\pi,\pi}$.

(b) $\lim_{m\to\infty}\int_{-\pi}^{\pi}|f-\sigma_m|=0$.

(c) Jika $g$ adalah fungsi yang dapat diintegralkan lainnya dengan koefisien Fourier
yang sama, maka $f\eae g$.

(d) Jika $x\in\ooint{-\pi,\pi}$ sedemikian rupa sehingga
$a=\lim_{t\in\dom f,t\uparrow
x}f(t)$ dan $b=\lim_{t\in\dom f,t\downarrow x}f(t)$ keduanya didefinisikan dalam
$\Bbb C$, maka

\Centerline{$\lim_{m\to\infty}\sigma_m(x)=\Bover12(a+b)$.}

(e) Jika $a=\lim_{t\in\dom f,t\uparrow \pi}f(t)$ dan $b=\lim_{t\in\dom
f,t\downarrow -\pi}f(t)$ keduanya didefinisikan dalam $\Bbb C$, maka

\Centerline{$\lim_{m\to\infty}\sigma_m(\pi)=\Bover12(a+b)$.}

(f) Jika $f$ terdefinisi dan kontinu di $x\in\ooint{-\pi,\pi}$, maka

\Centerline{$\lim_{m\to\infty}\sigma_m(x)=f(x)$.}

(g) Jika $\tilde f$, perpanjangan periodik $f$, terdefinisi dan
kontinu di $\pi$, maka

\Centerline{$\lim_{m\to\infty}\sigma_m(\pi)=f(\pi)$.}

\proof{{\bf (a)} Kita hanya perlu menggunakan 223D untuk melihat bahwa

$$\eqalign{\limsup_{\delta\downarrow 0}\Bover1{\delta}\int_0^{\delta}
&|f(x+t)+f(x-t)-2f(x)|dt\cr
&\le\limsup_{\delta\downarrow 0}
\Bover1{\delta}\bigl(\int_0^{\delta}|f(x+t)-f(x)|dt
+\int_0^{\delta}|f(x-t)-f(x)|dt\bigr)\cr
&=\limsup_{\delta\downarrow 0}
\Bover1{\delta}\int_{-\delta}^{\delta}|f(x+t)-f(x)|dt
=0\cr}$$

\noindent untuk hampir setiap $x\in\ooint{-\pi,\pi}$.

\medskip

{\bf (b)} Selanjutnya perhatikan bahwa, dalam bahasa 255O,

\Centerline{$\sigma_m=f*\psi_m$,}

\noindent dengan rumus terakhir di 282Db.   Akibatnya, oleh 255Od,

\Centerline{$\|\sigma_m\|_1\le\|f\|_1\|\psi_m\|_1$,}

\noindent menulis $\|\sigma_m\|_1=\int_{-\pi}^{\pi}|\sigma_m|$.
Tapi ini berarti kita punya

\Centerline{$f(x)
=\lim_{m\to\infty}\sigma_m(x)$ untuk hampir setiap $x$,
\quad$\limsup_{m\to\infty}\|\sigma_m\|_1\le\|f\|_1$;}

\noindent dan 245H mengikuti
$\lim_{m\to\infty}\|f-\sigma_m\|_1=0$.

\medskip

{\bf (c)} Jika $g$ mempunyai koefisien Fourier yang sama dengan $f$, maka
mempunyai jumlah Fourier dan Fej\'er yang sama, jadi kita punya

\Centerline{$g(x)=\lim_{m\to\infty}\sigma_m(x)=f(x)$}

\noindent hampir di mana-mana.

\medskip

{\bf (d)-(e)} Keduanya sama dengan mempertimbangkan $x\in\ocint{-\pi,\pi}$
sedemikian rupa sehingga

\Centerline{$\lim_{t\in\dom\tilde f,t\uparrow x}\tilde f(t)=a$,\quad
$\lim_{t\in\dom\tilde f,t\downarrow x}\tilde f(t)=b$.}

\noindent Dengan menetapkan $c={1\over 2}(a+b)$,
$\phi(t)=|\tilde f(x+t)+\tilde f(x-t)-2c|$ setiap kali ini didefinisikan, kami
memiliki $\lim_{t\in\dom\phi,t\downarrow 0}\phi(t)=0$, jadi pasti
$\lim_{\delta\downarrow 0}\bover1{\delta}\int_0^{\delta}\phi=0$, dan teorema
berlaku.

\medskip

{\bf (f)-(g)} adalah kasus khusus dari (d) dan (e).
}%end of proof of 282I

\leader{282J}{}\cmmnt{ Sekarang saya beralih ke kondisi konvergensi jumlah Fourier.    Mungkin hasil termudah -- yang
mencolok dan memuaskan -- adalah sebagai berikut.

\medskip

\noindent}{\bf Teorema} Misalkan $f$ adalah fungsi bernilai kompleks yang
terintegralkan-kuadrat pada $\ocint{-\pi,\pi}$. Misalkan
$\langle c_k\rangle_{k\in\Bbb Z}$ adalah koefisien Fourier-nya dan
$\sequencen{s_n}$ jumlah Fourier-nya\cmmnt{ (282A)}. Maka

(i) $\sum_{k=-\infty}^{\infty}|c_k|^2
=\Bover1{2\pi}\int_{-\pi}^{\pi}|f|^2$,

(ii) $\lim_{n\to\infty}\int_{-\pi}^{\pi}|f-s_n|^2=0$.

\proof{{\bf (a)} Saya ingat beberapa notasi dari 244N/244P.   Misalkan
$\eusm L_{\Bbb C}^2$
adalah ruang fungsi bernilai kompleks yang terintegralkan-kuadrat di
$\ocint{-\pi,\pi}$.   Untuk $g$, $h\in\eusm L_{\Bbb C}^2$, tulis

$$(g|h)=\int_{-\pi}^{\pi}g\times\bar h,\quad
\|g\|_2=\sqrt{(g|g)}.$$

\noindent Ingatlah bahwa $\|g+h\|_2\le\|g\|_2+\|h\|_2$ untuk semua $g$,
$h\in\eusm L_{\Bbb C}^2$ (244Fb/244Pb).   Untuk $k\in\Bbb Z$,
$x\in \ocint{-\pi,\pi}$ tetapkan $e_k(x)=e^{ikx}$, sehingga

$$(f|e_k)=\int_{-\pi}^{\pi}f(x)e^{-ikx}dx=2\pi c_k.$$

\noindent Selain itu, jika $|k|\le n$,

$$(s_n|e_k)=\sum_{j=-n}^nc_j\int_{-\pi}^{\pi}e^{ijx}e^{-ikx}dx
=2\pi c_k,$$

\noindent karena

$$\eqalign{\int_{-\pi}^{\pi}e^{ijx}e^{-ikx}dx&=2\pi\text{ jika }j=k,\cr
&=0\text{ jika }j\ne k.\cr}$$

\noindent  Jadi

\Centerline{$(f-s_n|e_k)=0$ setiap kali $|k|\le n$;}

\noindent khususnya,

$$(f-s_n|s_n)=\sum_{k=-n}^n\bar c_k(f-s_n|e_k)=0$$

\noindent untuk setiap $n\in\Bbb N$.

\medskip

{\bf (b)} Tetapkan $\epsilon>0$.   Unsur pembuktian selanjutnya adalah fakta
bahwa ada $m\in\Bbb N$, $a_{-m},\ldots,a_m\in\Bbb C$ sedemikian sehingga
$\|f-h\|_2\le\epsilon$, di mana $h=\sum_{k=-m}^ma_ke_k$.   \Prf\ Oleh
244Hb/244Pb kita tahu bahwa ada fungsi kontinu
$g:[-\pi,\pi]\to\Bbb C$ sehingga $\|f-g\|_2\le{{\epsilon}\over 3}$.
Selanjutnya, memodifikasi $g$ dalam interval pendek yang sesuai
$\ocint{\pi-\delta,\pi}$, kita dapat menemukan fungsi kontinu
$g_1:[-\pi,\pi]\to\Bbb C$ sehingga $\|g-g_1\|_2\le{{\epsilon}\over 3}$
dan $g_1(-\pi)=g_1(\pi)$.
(Atur $M=\sup_{x\in[-\pi,\pi]}|g(x)|$, ambil $\delta\in\ocint{0,2\pi}$
sedemikian rupa sehingga $(2M)^2\delta\le(\epsilon/3)^2$, dan atur
$g_1(\pi-t\delta)=tg(\pi-\delta)+(1-t)g(-\pi)$ untuk $t\in[0,1]$.)
Entah oleh
teorema Stone-Weierstrass (281J), atau menurut 282G di atas, ada
$a_{-m},\ldots,a_m$ seperti
 $|g_1(x)-\sum_{k=-m}^ma_ke^{ikx}|\le\Bover{\epsilon}{3\sqrt{2\pi}}$
untuk setiap $x\in[-\pi,\pi]$;  dengan menetapkan $h=\sum_{k=-m}^ma_ke_k$, kita punya
$\|g_1-h\|_2\le\bover13\epsilon$, sehingga

\Centerline{$\|f-h\|_2\le\|f-g\|_2+\|g-g_1\|_2+\|g_1-h\|_2
\le\epsilon$.   \Qed}

\medskip

{\bf (c)} Sekarang ambil sembarang $n\ge m$.   Maka $s_n-h$ merupakan kombinasi linear
dari  $e_{-n},\ldots,e_n$, jadi $(f-s_n|s_n-h)=0$.   Akibatnya

$$\eqalign{\epsilon^2
&\ge(f-h|f-h)\cr
&=(f-s_n|f-s_n)+(f-s_n|s_n-h)+(s_n-h|f-s_n)+(s_n-h|s_n-h)\cr
&=\|f-s_n\|_2^2+\|s_n-h\|_2^2
\ge\|f-s_n\|_2^2.\cr}$$

\noindent Jadi $\|f-s_n\|_2\le\epsilon$ untuk setiap $n\ge m$.   Karena
$\epsilon$ bersifat arbitrer, sehingga $\lim_{n\to\infty}\|f-s_n\|_2^2=0$, sebagaimana
dinyatakan dalam (ii).

\medskip

{\bf (d)} Adapun (i), kita punya

$$\sum_{k=-n}^n|c_k|^2=\Bover1{2\pi}\sum_{k=-n}^n\bar c_k(s_n|e_k)
=\Bover1{2\pi}(s_n|s_n)=\Bover1{2\pi}\|s_n\|_2^2.$$

\noindent Tapi tentu saja

\Centerline{$\bigl|\|s_n\|_2-\|f\|_2\bigr|\le\|s_n-f\|_2\to 0$}

\noindent sebagai $n\to\infty$, jadi

$$\sum_{k=-\infty}^{\infty}|c_k|^2
=\Bover1{2\pi}\lim_{n\to\infty}\|s_n\|_2^2=\Bover1{2\pi}\|f\|_2^2
=\Bover1{2\pi}\int_{-\pi}^{\pi}|f|^2,$$

\noindent sesuai kebutuhan.
}%end of proof of 282J

\leader{282K}{Akibat} Misalkan $L_{\Bbb C}^2$ adalah ruang Hilbert kelas-kelas
ekuivalensi fungsi bernilai kompleks yang terintegralkan-kuadrat pada
$\ocint{-\pi,\pi}$, dengan hasil kali dalam

$$(f^{\ssbullet}|g^{\ssbullet})
=\int_{-\pi}^{\pi}f\times\bar g$$

\noindent dan norma

$$\|f^{\ssbullet}\|_2
=\bigl(\int_{-\pi}^{\pi}|f|^2\bigr)^{1/2},$$

\noindent menulis $f^{\ssbullet}\in L_{\Bbb C}^2$ untuk kelas kesetaraan
dari fungsi terintegralkan-kuadrat
 $f$.   Misalkan $\ell_{\Bbb C}^2(\Bbb Z)$ menjadi
Ruang Hilbert dari barisan kompleks berujung ganda
yang dapat dijumlahkan persegi, dengan hasil kali dalam

$$(\pmb{c}|\pmb{d})=\sum_{k=-\infty}^{\infty}c_k\bar d_k$$

\noindent dan norma

$$\|\pmb{c}\|_2
=\bigl(\sum_{k=-\infty}^{\infty}|c_k|^2\bigr)^{1/2}$$

\noindent untuk $\pmb{c}=\langle c_k\rangle_{k\in\Bbb Z}$,
$\pmb{d}=\langle d_k\rangle_{k\in\Bbb Z}$ di $\ell_{\Bbb C}^2(\Bbb Z)$.
Maka kita mempunyai isomorfisme ruang hasil kali dalam
$S:L_{\Bbb C}^2\to\ell_{\Bbb C}^2(\Bbb Z)$
didefinisikan oleh

\Centerline{$S(f^{\ssbullet})(k)
=\Bover1{\sqrt{2\pi}}\int_{-\pi}^{\pi}f(x)e^{-ikx}dx$}

\noindent untuk setiap fungsi terintegralkan-kuadrat $f$ dan setiap
$k\in\Bbb Z$.

\proof{{\bf (a)} Seperti pada 282J, tulis $\eusm L_{\Bbb C}^2$ untuk ruang fungsi terintegralkan-kuadrat.   Jika $f$, $g\in\eusm L_{\Bbb C}^2$ dan
$f^{\ssbullet}=g^{\ssbullet}$, maka $f\eae g$, jadi

\Centerline{$\Bover1{\sqrt{2\pi}}\int_{-\pi}^{\pi}f(x)e^{-ikx}dx
=\Bover1{\sqrt{2\pi}}\int_{-\pi}^{\pi}g(x)e^{-ikx}dx$}

\noindent untuk setiap $k\in\Bbb Z$.    Jadi $S$ terdefinisi dengan baik.

\medskip

{\bf (b)}  $S$ linier.   \Prf\ Ini dasar.   Jika $f$,
$g\in\eusm L_{\Bbb C}^2$ dan $c\in\Bbb C$,

$$\eqalign{S(f^{\ssbullet}+g^{\ssbullet})(k)
&=\Bover1{\sqrt{2\pi}}\int_{-\pi}^{\pi}(f(x)+g(x))e^{-ikx}dx\cr
&=\Bover1{\sqrt{2\pi}}\int_{-\pi}^{\pi}f(x)e^{-ikx}dx
  +\Bover1{\sqrt{2\pi}}\int_{-\pi}^{\pi}g(x)e^{-ikx}dx\cr
&=S(f^{\ssbullet})(k)+S(g^{\ssbullet})(k)\cr}$$

\noindent untuk setiap $k\in\Bbb Z$, sehingga
$S(f^{\ssbullet}+g^{\ssbullet})=S(f^{\ssbullet})+S(g^{\ssbullet})$.
Demikian pula,

\Centerline{$S(cf^{\ssbullet})(k)
=\Bover1{\sqrt{2\pi}}\int_{-\pi}^{\pi}cf(x)e^{-ikx}dx
=\Bover{c}{\sqrt{2\pi}}\int_{-\pi}^{\pi}f(x)e^{-ikx}dx
=cS(f^{\ssbullet})(k)$}

\noindent untuk setiap $k\in\Bbb Z$, sehingga
$S(cf^{\ssbullet})=cS(f^{\ssbullet})$.   \Qed

\medskip

{\bf (c)} Jika $f\in\eusm L_{\Bbb C}^2$ mempunyai koefisien Fourier $c_k$,
maka
$S(f^{\ssbullet})=\langle c_k\sqrt{ 2\pi}\rangle_{k\in\Bbb Z}$, maka oleh 282J(i)

$$\|S(f^{\ssbullet})\|^2_2
=2\pi\sum_{k=-\infty}^{\infty}|c_k|^2
=\int_{-\pi}^{\pi}|f|^2
=\|f^{\ssbullet}\|_2^2.$$

\noindent Jadi $Su\in\ell_{\Bbb C}^2(\Bbb Z)$ dan $\|Su\|_2=\|u\|_2$ untuk
setiap $u\in L_{\Bbb C}^2$.   Karena $S$ bersifat linier dan menjaga norma,
pastinya bersifat injektif.

\medskip

{\bf (d)}  Sekarang $(Sv|Su)=(v|u)$ untuk setiap $u$, $v\in
L_{\Bbb C}^2$.   \Prf\ (Ini tentu saja merupakan fakta standar tentang ruang Hilbert.) Kita tahu bahwa untuk setiap $t\in\Bbb R$

$$\eqalign{\|u\|_2^2+2\Real(e^{it}(v|u))+\|v\|_2^2
&=(u|u)+e^{it}(v|u)+e^{-it}(u|v)+(v|v)\cr
&=(u+e^{it}v|u+e^{it}v)\cr
&=\|u+e^{it}v\|_2^2
=\|S(u+e^{it}v)\|_2^2\cr
&=\|Su\|_2^2+2\Real(e^{it}(Sv|Su))+\|Sv\|_2^2\cr
&=\|u\|_2^2+2\Real(e^{it}(Sv|Su))+\|v\|_2^2,\cr}$$

\noindent sehingga $\Real(e^{it}(Sv|Su))=\Real(e^{it}(v|u))$.   Karena $t$ adalah
sewenang-wenang, $(Sv|Su)=(v|u)$.   \Qed

\medskip

{\bf (e)}    Terakhir, $S$ bersifat surjektif.   \Prf\ Misalkan
$\pmb{c}=\langle c_k\rangle_{k\in\Bbb Z}$ menjadi anggota
$\ell_{\Bbb C}^2(\Bbb Z)$.   Tetapkan
$c^{(n)}_k=c_k$ jika $|k|\le n$, $0$ sebaliknya, dan
$\pmb{c}^{(n)}=\langle c^{(n)}_k\rangle_{k\in\Bbb Z}$.   Mempertimbangkan

$$s_n=\sum_{k=-n}^nc_ke_k,\quad u_n=s_n^{\ssbullet}$$

\noindent tempat saya menulis $e_k(x)=\bover1{\sqrt{2\pi}}e^{ikx}$ untuk
$x\in\ocint{-\pi,\pi}$.   Maka $Su_n=\pmb{c}^{(n)}$, dengan perhitungan
yang sama seperti pada bagian (a) pembuktian 282J.   Sekarang

\Centerline{$\|\pmb{c}^{(n)}-\pmb{c}\|_2=\sqrt{\sumop_{|k|>n}|c_k|^2}
\to 0$}

\noindent sebagai $n\to\infty$, jadi

\Centerline{$\|u_m-u_n\|_2=\|\pmb{c}^{(m)}-\pmb{c}^{(n)}\|_2
\to 0$}

\noindent sebagai $m$, $n\to\infty$, dan $\sequencen{u_n}$ adalah barisan Cauchy
di $L_{\Bbb C}^2$.   Karena $L_{\Bbb C}^2$ sudah lengkap
(244G/244Pb), $\sequencen{u_n}$ mempunyai batas $u\in L_{\Bbb C}^2$, dan sekarang

\Centerline{$Su=\lim_{n\to\infty}Su_n=\lim_{n\to\infty}\pmb{c}^{(n)}
=\pmb{c}$.   \Qed}

Jadi $S:L_{\Bbb C}^2\to\ell_{\Bbb C}^2(\Bbb Z)$ adalah isomorfisme ruang hasil kali dalam.
}%end of proof of 282K

\cmmnt{\medskip

\noindent{\bf Catatan} Dalam bahasa ruang Hilbert, semua yang
yang terjadi di sini adalah bahwa $\langle e_k^{\ssbullet}\rangle_{k\in\Bbb Z}$ adalah
sebuah `basis ruang Hilbert' atau `barisan ortonormal lengkap' di
$L_{\Bbb C}^2$, yang dicocokkan oleh $S$ dengan basis standar
$\ell^2_{\Bbb C}(\Bbb Z)$.   Satu-satunya langkah yang memerlukan analisis real
non-sepele, karena
bertentangan dengan teori umum ruang Hilbert, adalah pemeriksaan bahwa subruang linier
yang dihasilkan oleh $\{e_k^{\ssbullet}:k\in\Bbb Z\}$ padat;
ini adalah bagian (b) dari bukti 282J.

Perhatikan bahwa meskipun $S:L^2\to\ell^2$ mudah dijelaskan, kebalikannya adalah
yang lebih merupakan masalah.   Jika $\pmb{c}\in\ell^2$, kami ingin mengatakan bahwa
$S^{-1}\pmb{c}$ adalah kelas kesetaraan $f$, di mana
$f(x)=\bover1{\sqrt{2\pi}}\sum_{k=-\infty}^{\infty}c_ke^{ikx}$ untuk setiap
$x$.   Ini berfungsi dengan baik jika $\{k:c_k\ne 0\}$ terbatas, tetapi untuk kasus umum
kurang jelas bagaimana menafsirkan jumlahnya.   Faktanya adalah
jika $\pmb{c}\in\ell^2$ maka

\Centerline{$g(x)
=\Bover1{\sqrt{2\pi}}\lim_{n\to\infty}\sum_{k=-n}^nc_ke^{ikx}$}

\noindent didefinisikan untuk hampir setiap $x\in\ocint{-\pi,\pi}$, dan
$S^{-1}\pmb{c}=g^{\ssbullet}$ di $L^2$;  ini sebenarnya adalah teorema
Carleson (286V).   Bukti teorema Carleson berada di luar jangkauan kami untuk
 saat ini.   Hasil yang tercakup dalam bagian ini adalah bahwa

\Centerline{$h(x)
=\Bover1{\sqrt{2\pi}}\lim_{m\to\infty}\Bover1{m+1}\sum_{n=0}^m
\sum_{k=-n}^nc_ke^{ikx}$}

\noindent didefinisikan untuk hampir setiap $x\in\ocint{-\pi,\pi}$, dan
$h^{\ssbullet}=S^{-1}\pmb{c}$.   (Pada dasarnya, hasil yang baru saja dibuktikan
menunjukkan bahwa ada {\it suatu} fungsi $f$ yang terintegralkan-kuadrat sedemikian sehingga
$\pmb{c}$ adalah barisan koefisien Fourier dari $f$; sekarang 282Ia menyatakan
bahwa jumlah Fej\'er dari $f$ konvergen hampir di mana-mana ke $f$, yaitu
$h\eae\bover1{\sqrt{2\pi}}f$.)
}%end of comment

\leader{282L}{}\cmmnt{ Hasil berikutnya adalah yang paling mudah, dan salah satu teorema
yang paling berguna, mengenai konvergensi titik demi titik dari jumlah Fourier.

\medskip

\noindent}{\bf Teorema} Misalkan $f$ adalah fungsi bernilai kompleks yang
dapat diintegralkan pada $\ocint{-\pi,\pi}$, dan $\sequencen{s_n}$ barisan
jumlah Fourier-nya.

(i) Jika $f$ terdiferensiasi di $x\in\ooint{-\pi,\pi}$, maka
$f(x)=\lim_{n\to\infty}s_n(x)$.

(ii) Jika perpanjangan periodik $\tilde f$ dari $f$ dapat terdiferensiasi di
$\pi$, lalu $f(\pi)=\lim_{n\to\infty}s_n(\pi)$.

\wheader{282L}{0}{0}{0}{60pt}

\proof{{\bf (a)} Ambil $x\in\ocint{-\pi,\pi}$ sedemikian rupa sehingga $\tilde
f$ terdiferensiasi di $x$;  tentu saja ini mencakup kedua bagian.    Kami
punya

$$s_n(x)=\Bover1{2\pi}\int_{-\pi}^{\pi}
\Bover{\tilde f(x+t)}{\sin{1\over 2}t}\sin(n+\Bover12)t\,dt$$

\noindent untuk setiap $n$, oleh 282Da.

\medskip

{\bf (b)} Selanjutnya,

$$\int_{-\pi}^{\pi}\Bover{\tilde f(x+t)-\tilde f(x)}{t}dt$$

\noindent ada di $\Bbb C$, karena pasti ada beberapa
$\delta\in\ocint{0,\pi}$
sehingga $(\tilde f(x+t)-\tilde f(x))/t$ dibatasi pada
$\{t:0<|t|\le\delta\}$, sedangkan

$$\int_{-\pi}^{-\delta}
\Bover{\tilde f(x+t)-\tilde f(x)}{t}dt,
\quad\int_{\delta}^{\pi}\Bover{\tilde f(x+t)-\tilde f(x)}{t}dt$$

\noindent ada karena $1/t$ dibatasi pada interval tersebut.
akibatnya

$$\int_{-\pi}^{\pi}
\Bover{\tilde f(x+t)-\tilde f(x)}{\sin{1\over 2}t}dt$$

\noindent ada, karena $|t|\le\pi|\sin{1\over 2}t|$ jika
$|t|\le\pi$.   Jadi dengan lema Riemann-Lebesgue (282Fb),

$$\lim_{n\to\infty}\int_{-\pi}^{\pi}
\Bover{\tilde f(x+t)-\tilde f(x)}{\sin{1\over 2}t}\sin(n+\Bover12)t\,dt
=0.$$

\medskip

{\bf (c)} Karena

$$\Bover1{2\pi}\int_{-\pi}^{\pi}\tilde f(x)
  \Bover{\sin(n+{1\over 2}\pushbottom{3.5pt})t}{\sin{1\over 2}t}dt
=\tilde f(x)$$

\noindent untuk setiap $n$ (282Dc),

$$s_n(x)
=\tilde f(x)+\Bover{1}{2\pi}\int_{-\pi}^{\pi}
  \Bover{\tilde f(x+t)-\tilde f(x)}{\sin{1\over 2}t}
  \sin(n+\Bover12)t\,dt
\to\tilde f(x)$$

\noindent sebagai $n\to\infty$, sesuai kebutuhan.
}%end of proof of 282L

\leader{282M}{Lema} Misalkan $f$ adalah fungsi
bernilai kompleks, didefinisikan hampir di mana-mana dan variasi terbatas pada
$\ocint{-\pi,\pi}$.   Maka $\sup_{k\in\Bbb Z}|kc_k|<\infty$, di mana $c_k$
adalah $k$koefisien Fourier th dari $f$, seperti pada 282A.

\proof{ Tetapkan

\Centerline{$M
=\lim_{x\in\dom f,x\uparrow \pi}|f(x)|+\Var_{\ooint{-\pi,\pi}}(f)$.}

\noindent Oleh 224J,

$$\eqalign{|kc_k|
&=\Bover{1}{2\pi}\bigl|\int_{-\pi}^{\pi}kf(t)e^{-ikt}dt\bigr|
\le\Bover1{2\pi}M
   \sup_{c\in[-\pi,\pi]}\bigl|\int_{-\pi}^cke^{-ikt}dt\bigr|\cr
&=\Bover{M}{2\pi}\sup_{c\in[-\pi,\pi]}|e^{-ikc}-e^{ik\pi}|
\le\Bover{M}{\pi}\cr}$$

\noindent untuk setiap $k$.
}%end of proof of 282M

\leader{282N}{}\cmmnt{ Saya memberikan lema lain, mengekstraksi bagian teknis
dari
bukti teorema berikutnya.   (Aplikasi paling alami ada di
282Xn.)

\medskip

\noindent}{\bf Lema} Misalkan $\sequence{k}{d_k}$ adalah barisan bilangan kompleks, dan
tetapkan $t_n=\sum_{k=0}^nd_k$, $\tau_m=\bover1{m+1}\sum_{n=0}^mt_n$ untuk $n$,
$m\in\Bbb N$.   Misalkan $\sup_{k\in\Bbb N}|kd_k|=M<\infty$.   Maka
untuk $j\ge 1$ dan setiap $c\in\Bbb C$,

\Centerline{$|t_n-c|\le\Bover{M}{j}+(2j+3)\sup_{m\ge n-n/j}|\tau_m-c|$}

\noindent untuk setiap $n\ge j^2$.

\proof{{\bf (a)} Hal pertama yang perlu diperhatikan adalah untuk $n$,
$n'\in\Bbb N$,

\Centerline{$|t_n-t_{n'}|\le\Bover{M|n-n'|}{1+\min(n,n')}$.}

\noindent\Prf\ Jika $n=n'$ ini sepele.   Misalkan $n'<n$.   Maka

$$|t_n-t_{n'}|
=|\sum_{k=n'+1}^{n}d_k|
\le\sum_{k=n'+1}^n\Bover{M}{k}
\le\Bover{M(n-n')}{n'+1}
=\Bover{M|n-n'|}{1+\min(n',n)}.$$

\noindent Tentu saja kasus $n<n'$ identik.   \Qed


\medskip

{\bf (b)} Sekarang ambil sembarang $n\ge j^2$.
Tetapkan $\eta=\sup_{m\ge n-n/j}|\tau_m-c|$.   Misalkan $m\ge j$ menjadi
$jm\le n<j(m+1)$;  lalu $n<jm+m$;  Juga

\Centerline{$n(1-\bover1{j})\le m(j+1)(1-{1\over j})\le mj$.}

\noindent    Tetapkan

$$\tau^*=\Bover1m\sum_{n'=jm+1}^{jm+m}t_{n'}
=\Bover{jm+m+1}{m}\tau_{jm+m}-\Bover{jm+1}{m}\tau_{jm}.$$

\noindent Lalu

$$\eqalign{|\tau^*-c|
&=|\Bover{jm+m+1}{m}\tau_{jm+m}-\Bover{jm+1}{m}\tau_{jm}-c|\cr
&=|\Bover{jm+m+1}{m}(\tau_{jm+m}-c)
   -\Bover{jm+1}{m}(\tau_{jm}-c)|\cr
&\le\Bover{jm+m+1}{m}\eta+\Bover{jm+1}{m}\eta
\le(2j+3)\eta.\cr}$$

\noindent Di sisi lain,

$$\eqalign{|\tau^*-t_n|
&=\bigl|\Bover1m\sum_{n'=jm+1}^{jm+m}(t_{n'}-t_n)\bigr|
\le\Bover1m\sum_{n'=jm+1}^{jm+m}\Bover{M|n-n'|}{1+\min(n,n')}\cr
&\le\Bover1m\sum_{n'=jm+1}^{jm+m}\Bover{Mm}{1+jm}
=\Bover{Mm}{1+jm}
\le\Bover{M}{j}.\cr}$$

\noindent Dengan menggabungkan ini, kita punya

\Centerline{$|t_n-c|\le|t_n-\tau^*|+|\tau^*-c|
\le\Bover{M}{j}+(2j+3)\eta
=\Bover{M}{j}+(2j+3)\sup_{m\ge n-n/j}|\tau_m-c|$,}

\noindent sesuai kebutuhan.
}%end of proof of 282N

\leader{282O}{Teorema} Misalkan $f$ adalah fungsi bernilai kompleks dengan
variasi terbatas, terdefinisi hampir di mana-mana pada $\ocint{-\pi,\pi}$,
dan misalkan $\sequencen{s_n}$ adalah barisan jumlah Fourier-nya.

(i) Jika $x\in\ooint{-\pi,\pi}$, maka

\Centerline{$\lim_{n\to\infty}s_n(x)
=\Bover12(\lim_{t\in\dom f,t\uparrow x}f(t)
  +\lim_{t\in\dom f,t\downarrow x}f(t))$.}

(ii) $\lim_{n\to\infty}s_n(\pi)
=\Bover12(\lim_{t\in\dom f,t\uparrow\pi}f(t)
  +\lim_{t\in\dom f,t\downarrow-\pi}f(t))$.

(iii) Jika $f$ didefinisikan di seluruh $\ocint{-\pi,\pi}$, kontinu,
dan $\lim_{t\downarrow -\pi}f(t)=f(\pi)$, maka $s_n(x)\to f(x)$
seragam di $\ocint{-\pi,\pi}$.

\proof{{\bf (a)} Perhatikan terlebih dahulu bahwa 224F menunjukkan bahwa limit
$\lim_{t\in\dom f,t\downarrow x}f(t)$,
$\lim_{t\in\dom f,t\uparrow x}f(t)$ yang diperlukan dalam rumus di atas
selalu ada.   Kita juga mengetahui dari 282M bahwa $M=
\penalty-100
\sup_{k\in\Bbb Z}|kc_k|<\infty$, di mana $c_k$ adalah koefisien Fourier ke-$k$
dari $f$.

Ambil sembarang $x\in\ocint{-\pi,\pi}$, dan tetapkan

\Centerline{$c=\bover12(\lim_{t\in\dom f,t\uparrow x}\tilde
f(t)+\lim_{t\in\dom\tilde f,t\downarrow x}\tilde f(t))$,}

\noindent menulis $\tilde f$ untuk perpanjangan berkala $f$, sebagai
biasa.   Menurut 282Id-282Ie,
$c=\lim_{m\to\infty}\sigma_m(x)$, menulis $\sigma_m$ untuk Fej\'er
jumlah $f$.   Misalkan $\epsilon>0$.   Ambil sembarang $j\ge\max(2,2M/\epsilon)$, dan
$m_0\ge 1$ sedemikian sehingga
$|\sigma_m(x)-c|\le\epsilon/(2j+3)$ untuk setiap $m\ge m_0$.

Sekarang jika $n\ge\max(j^2,2m_0)$, terapkan Lema 282N dengan

\Centerline{$d_0=c_0$,
\quad$d_k=c_ke^{ikx}+c_{-k}e^{-ikx}$ untuk $k\ge 1$,}

\noindent sehingga $t_n=s_n(x)$, $\tau_m=\sigma_m(x)$ dan $|kd_k|\le 2M$
untuk setiap $k$, $n$, $m\in\Bbb N$.   Kita mempunyai $n-n/j\ge{1\over 2}n\ge
m_0$, jadi

\Centerline{$\eta=\sup_{m\ge n-n/j}|\tau_m-c|
\le\sup_{m\ge m_0}|\tau_m-c|\le\Bover{\epsilon}{2j+3}$.}

\noindent Jadi, menurut 282N,

\Centerline{$|s_n(x)-c|=|t_n-c|\le
\Bover{2M}{j}+(2j+3)\sup_{m\ge n-n/j}|\tau_m-c|\le\epsilon+(2j+3)\eta
\le 2\epsilon$.}

\noindent Karena $\epsilon$ bersifat arbitrer, $\lim_{n\to\infty}s_n(x)=c$, sesuai kebutuhan.

\medskip

{\bf (b)} Ini membuktikan (i) dan (ii) teorema ini.
Akhirnya, untuk (iii), amati bahwa dalam kondisi ini $\sigma_m(x)\to
f(x)$ seragam sebagai $m\to\infty$, oleh 282G.   Jadi dengan $\epsilon>0$
kita memilih $j\ge \max(2,2M/\epsilon)$ dan $m_0\in\Bbb N$ seperti
sehingga $|\sigma_m(x)-f(x)|\le\epsilon/(2j+3)$ setiap kali $m\ge m_0$ dan
$x\in\ocint{-\pi,\pi}$.   Dengan perhitungan yang sama seperti sebelumnya,

\Centerline{$|s_n(x)-f(x)|\le 2\epsilon$}

\noindent untuk setiap $n\ge\max(j^2,2m_0)$ dan setiap
$x\in\ocint{-\pi,\pi}$.   Karena $\epsilon$ bersifat arbitrer,
$\lim_{n\to\infty}s_n(x)=f(x)$ seragam untuk $x\in\ocint{-\pi,\pi}$.
}%end of proof of 282O

\leader{282P}{Akibat} Misalkan $f$ adalah fungsi bernilai kompleks yang dapat
diintegralkan pada $\ocint{-\pi,\pi}$, dan $\sequencen{s_n}$ barisan jumlah
Fourier-nya.

(i) Misalkan $x\in\ooint{-\pi,\pi}$ sedemikian rupa sehingga $f$ merupakan variasi
yang dibatasi pada beberapa lingkungan $x$.   Maka

\Centerline{$\lim_{n\to\infty}s_n(x)=\Bover12(\lim_{t\in\dom f,t\uparrow
x}f(t)+\lim_{t\in\dom f,t\downarrow x}f(t))$.}

(ii) Jika terdapat $\delta>0$ sehingga $f$ memiliki variasi terbatas pada
baik $\ocint{-\pi,-\pi+\delta}$ maupun $[\pi-\delta,\pi]$, maka

\Centerline{$\lim_{n\to\infty}s_n(\pi)
=\Bover12(\lim_{t\in\dom f,t\uparrow \pi}f(t)
  +\lim_{t\in\dom f,t\downarrow-\pi}f(t))$.}

\proof{ Dalam kasus (i), ambil $\delta>0$ sehingga $f$ merupakan variasi terbatas
pada $[x-\delta,x+delta]$ dan tetapkan $f_1(t)=f(t)$ jika
$t\in\dom f\cap[x-\delta,x+\delta]$, $0$ untuk
$t\in\ocint{-\pi,\pi}$ lainnya; dalam kasus (ii), tetapkan $f_1(t)=f(t)$ jika $t\in\dom
f$ dan $|t|\ge\pi-\delta$, serta $0$ untuk $t\in\ocint{-\pi,\pi}$ lainnya, lalu tetapkan
bahwa $x=\pi$.   Dalam kedua kasus tersebut, $f_1$ memiliki variasi terbatas, sehingga dengan
Menurut 282O, jumlah Fourier $\sequencen{s'_n}$ dari $f_1$ konvergen di
$x$ ke nilai yang diberikan oleh rumus di atas.   Namun sekarang amati bahwa,
menulis $\tilde f$ dan $\tilde f_1$ untuk ekstensi periodik $f$
dan $f_1$, $\tilde f-\tilde f_1=0$ di lingkungan $x$, jadi

$$\int_{-\pi}^{\pi}
  \Bover{\tilde f(x+t)-\tilde f_1(x+t)}{\sin{1\over 2}t}dt$$

\noindent ada di $\Bbb C$, dan oleh 282Fb

$$\lim_{n\to\infty}\int_{-\pi}^{\pi}
 \bover{\tilde f(x+t)-\tilde f_1(x+t)}{\sin{1\over 2}t}
\sin(n+\Bover12)t\,dt=0,$$

\noindent yaitu $\lim_{n\to\infty}s_n(x)-s'_n(x)=0$.   Jadi
$\sequencen{s_n}$ juga konvergen ke batas kanan.
}%end of proof of 282P

\leader{282Q}{}\cmmnt{ Saya tidak dapat meninggalkan bagian ini tanpa menyebutkan
salah satu fakta terpenting tentang deret Fourier, meskipun
tidak punya ruang di sini untuk membahas konsekuensinya.

\medskip

\noindent}{\bf Teorema} Misalkan $f$ dan $g$ adalah fungsi
bernilai kompleks yang dapat diintegralkan pada $\ocint{-\pi,\pi}$, dan
$\langle c_k\rangle_{k\in\Bbb Z}$,
$\langle d_k\rangle_{k\in\Bbb Z}$ koefisien Fourier-nya.   Misalkan
$f*g$ menjadi konvolusinya, yang ditentukan oleh rumus

$$(f*g)(x)=\int_{-\pi}^{\pi}f(x-_{2\pi}t)g(t)dt
=\int_{-\pi}^{\pi}\tilde f(x-t)g(t)dt,$$

\noindent\cmmnt{seperti pada 255O,} menulis $\tilde f$ untuk perpanjangan periodik
dari $f$.   Maka koefisien Fourier dari $f*g$
adalah $\langle 2\pi c_kd_k\rangle_{k\in\Bbb Z}$.

\proof{ Oleh 255O(c-i),

$$\eqalign{\Bover1{2\pi}\int_{-\pi}^{\pi}(f*g)(x)e^{-ikx}dx
&=\Bover1{2\pi}\int_{-\pi}^{\pi}\int_{-\pi}^{\pi}
e^{-ik(t+u)}f(t)g(u)dtdu\cr
&=\Bover1{2\pi}\int_{-\pi}^{\pi}e^{-ikt}f(t)dt\int_{-\pi}^{\pi}
e^{-iku}g(u)du
=2\pi c_kd_k.\cr}$$
}%end of proof of 282Q

\leader{*282R}{}\cmmnt{ Karena terburu-buru untuk membahas teorema konvergensi
dari jumlah Fej\'er dan Fourier, saya agak mengabaikan manipulasi dasar
yang penting saat menerapkan teori.
Salah satu hasil dasarnya adalah sebagai berikut.

\medskip

\noindent}{\bf Proposisi} (a) Misalkan $f:[-\pi,\pi]\to\Bbb C$ adalah fungsi
yang kontinu mutlak sehingga $f(-\pi)=f(\pi)$, dan
$\family{k}{\Bbb Z}{c_k}$ merupakan barisan koefisien Fourier-nya.   Maka
koefisien Fourier dari $f'$ adalah $\family{k}{\Bbb Z}{ikc_k}$.

(b)\discrversionA{\footnote{Dikoreksi dan disempurnakan 2009.}}{}
Misalkan $f:\Bbb R\to\Bbb C$ adalah fungsi terdiferensiasi sehingga $f'$
kontinu mutlak pada $[-\pi,\pi]$, dan $f(\pi)=f(-\pi)$.
Jika $\family{k}{\Bbb Z}{c_k}$ adalah koefisien Fourier dari
$f\restr\ocint{-\pi,\pi}$, maka
$\sum_{k=-\infty}^{\infty}|c_k|$ berhingga.

\proof{{\bf (a)} Oleh 225Cb, $f'$ dapat diintegralkan pada $[-\pi,\pi]$;  oleh
225E, $f$ merupakan integral tak tentu dari $f'$.   Jadi, menurut 225F,

\Centerline{$\Bover1{2\pi}\int_{-\pi}^{\pi}f'(x)e^{-ikx}dx
=\Bover{f(\pi)e^{-ik\pi}-f(-\pi)e^{ik\pi}}{2\pi}+ikc_k
=ikc_k$}

\noindent untuk setiap $k\in\Bbb Z$.

\medskip

{\bf (b)(i)} Misalkan dulu $f'(\pi)=f'(-\pi)$.
Dengan (a), diterapkan dua kali, koefisien Fourier $f''$ adalah
$\family{k}{\Bbb Z}{-k^2c_k}$, jadi $\sup_{k\in\Bbb Z}k^2|c_k|$ terbatas;
karena $\sum_{k=1}^{\infty}\Bover1{k^2}<\infty$,
$\sum_{k=-\infty}^{\infty}|c_k|<\infty$.

\medskip

\quad{\bf (ii)} Selanjutnya, anggaplah $f(x)=x^2$ untuk setiap $x$.   Lalu, untuk
$k\ne 0$,

$$\eqalign{c_k
&=\Bover1{2\pi}\int x^2e^{-ikx}dx
=\Bover1{2\pi}\bigl(-\Bover1{ik}(\pi^2e^{-ik\pi}-\pi^2e^{ik\pi})
   +\int_{-\pi}^{\pi}\Bover{2x}{ik}e^{-ikx}dx\bigr)\cr
&=\Bover1{ik\pi}\bigl(-\Bover1{ik}(\pi e^{-ik\pi}+\pi e^{ik\pi})
   +\Bover1{ik}\int_{-\pi}^{\pi}e^{ikx}dx\bigr)
=\Bover{2}{k^2}(-1)^k,\cr}$$

\noindent jadi
$\sum_{k\in\Bbb Z}|c_k|\le|c_0|+4\sum_{k=1}^{\infty}\Bover1{k^2}$
terbatas.

\medskip

\quad{\bf (iii)} Secara umum, kita dapat menyatakan $f$ sebagai $f_1+cf_2$ di mana
$f_2(x)=x^2$ untuk setiap $x$, $c=\Bover1{4\pi}(f'(\pi)-f'(-\pi))$, dan
$f_1$ memenuhi kondisi (i);  sehingga
$\family{k}{\Bbb Z}{c_k}$ adalah jumlah dari dua barisan yang dapat dijumlahkan dan
itu sendiri dapat dijumlahkan.
}%end of proof of 282R

\exercises{
\leader{282X}{Latihan dasar $\pmb{>}$(a)}
%\spheader 282Xa
Misalkan $\langle c_k\rangle_{k\in\Bbb Z}$
adalah barisan bilangan kompleks berujung ganda yang terjumlah mutlak.
Tunjukkan bahwa $f(x)=\sum_{k=-\infty}^{\infty}c_ke^{ikx}$ ada untuk setiap
$x\in\Bbb R$, bahwa $f$ kontinu dan periodik, dan koefisien Fourier
-nya adalah $c_k$.
%282A

\spheader 282Xc Tetapkan $\phi_n(t)=\bover2t\sin(n+\bover12t)$ untuk $t\ne 0$.
(Kadang-kadang ini disebut kernel Dirichlet {\bf yang dimodifikasi}.)
Tunjukkan bahwa untuk fungsi apa pun yang dapat diintegralkan $f$ pada $\ocint{-\pi,\pi}$, dengan
jumlah Fourier $\sequencen{s_n}$ dan ekstensi periodik $\tilde f$,

\Centerline{$\lim_{n\to\infty}|s_n(x)-\bover1{2\pi}
  \int_{-\pi}^{\pi}\phi_n(t)\tilde f(x+t)dt|
=0$}

\noindent untuk setiap $x\in\ocint{-\pi,\pi}$.   \Hint{menunjukkan bahwa
$\bover2t-\bover{1}{\sin{1\over 2}t}$ dibatasi, dan menggunakan 282E.}
%282E

\spheader 282Xd Berikan bukti 282Ib dari 242O, 255O dan 282G.
%282I

\spheader 282Xe Berikan bukti lain dari 282Ic, berdasarkan 222D, 281J dan ide
dalam bukti 242O dan bukan pada 282H.
%282I

\spheader 282Xf Gunakan gagasan 255Ya untuk mempersingkat salah satu langkah dalam bukti
dari 282H, dengan mengambil
\discrcenter{468pt}{$g_m(t)
=\min(\bover{m+1}{2\pi},\bover{\pi}{4(m+1)t^2})$ }untuk $|t|\le\delta$, sehingga
menjadi $g_m\ge\psi_m$ di $[-\delta,\delta]$.
%282H

\sqheader 282Xg(i) Misalkan $f$ adalah fungsi bernilai real yang terintegralkan-kuadrat
pada $\ocint{-\pi,\pi}$, dan $\sequence{k}{a_k}$,
$\langle b_k\rangle_{k\ge 1}$ merupakan koefisien Fourier real (282Ba).
Tunjukkan bahwa
$\bover12a_0^2+\sum_{k=1}^{\infty}(a_k^2+b_k^2)
=\bover1{\pi}\int_{-\pi}^{\pi}|f|^2$.  (ii) Tunjukkan bahwa
$f\mapsto(\sqrt{\bover{\pi}2}a_0,\sqrt{\pi}a_1,
\ifdim\pagewidth>467pt\penalty-100\fi
\sqrt{\pi}b_1,\ldots)$
mendefinisikan isomorfisme ruang hasil kali dalam antara ruang Hilbert
real $L^2_{\Bbb R}$ dari kelas ekivalensi fungsi
yang terintegralkan-kuadrat real di $\ocint{-\pi,\pi}$ dan ruang Hilbert real
$\ell^2_{\Bbb R}$ dari barisan yang dapat dijumlahkan kuadratnya.
%282K

\spheader 282Xh Tunjukkan bahwa
$\bover{\pi}4=1-\bover13+\bover15-\bover17+\ldots$.   \Hint{temukan deret Fourier
dari $f$ di mana $f(x)=x/|x|$, dan hitung jumlah deret
di $\bover{\pi}2$.   Tentu saja ada metode lain, misalnya
memeriksa deret Taylor untuk $\arctan 1=\bover{\pi}4$.}
%282L

\spheader 282Xi Misalkan $f$ adalah fungsi bernilai kompleks yang dapat diintegralkan
pada $\ocint{-\pi,\pi}$, dan $\sequencen{s_n}$ merupakan barisan jumlah Fourier.   Misalkan $x\in\ooint{-\pi,\pi}$, $a\in\Bbb C$ sedemikian rupa sehingga
$\biggerint_{-\pi}^{\pi}\Bover{f(t)-a}{t-x}dt$ ada dan terbatas.
Tunjukkan $\lim_{n\to\infty}s_n(x)=a$itu.
Jelaskan bagaimana hal ini menggeneralisasi 282L.   Modifikasi apa yang sesuai agar
mendapatkan batas $\lim_{n\to\infty}s_n(\pi)$?
%282L

\spheader 282Xj Misalkan $\alpha>0$, $K\ge 0$ dan
$f:\ooint{-\pi,\pi}\to\Bbb C$ sedemikian rupa sehingga
$|f(x)-f(y)|\le K|x-y|^{\alpha}$ untuk semua $x$, $y\in\ooint{-\pi,\pi}$.
(Fungsi tersebut disebut {\bf H\"older}kontinu.)
Tunjukkan bahwa jumlah Fourier dari $f$ konvergen ke $f$ di mana saja di
$\ooint{-\pi,\pi}$.   \Hint{gunakan 282Xi.}  (Bandingkan 282Yb.)
%282Xi 282L

\spheader 282Xk Dalam 282L, tunjukkan bahwa cukup jika $\tilde f$
terdiferensiasi sehubungan dengan domainnya di $x$ atau $\pi$ (lihat 262Fb),
daripada terdiferensiasi dalam arti sempit.
%282L

\spheader 282Xl Tunjukkan bahwa $\lim_{a\to\infty}\int_0^a\bover{\sin t}tdt$
ada dan terbatas.   \Hint{menggunakan 224J untuk memperkirakan
$\int_a^b\bover{\sin t}tdt$ untuk $0<a\le b$.}
%

\spheader 282Xm Tunjukkan bahwa
$\int_0^{\infty}\bover{|\sin t|}tdt=\infty$.   \Hint{menunjukkan bahwa
$\sup_{a\ge 0}|\int_1^a\bover{\cos 2t}tdt|<\infty$, dan oleh karena itu
$\sup_{a\ge 0}\int_1^a\bover{\sin^2t}{t}dt\penalty-100=\infty$.}
%282Xl

\sqheader 282Xn Misalkan $\sequence{k}{d_k}$ menjadi urutan dalam $\Bbb C$ seperti
sehingga $\sup_{k\in\Bbb N}|kd_k|<\infty$ dan

\Centerline{$\lim_{m\to\infty}\Bover1{m+1}\sum_{n=0}^m\sum_{k=0}^nd_k
=c\in\Bbb C$.}

\noindent Tunjukkan bahwa $c=\sum_{k=0}^{\infty}d_k$.   \Hint{282N.}
%282N

\sqheader 282Xo Tunjukkan bahwa
$\sum_{n=1}^{\infty}\bover1{n^2}=\bover{\pi^2}6$.   \Hint{(b-ii) dari bukti
dari 282R.}
%282Q

\spheader 282Xp Misalkan $f$ adalah fungsi bernilai kompleks
yang dapat diintegralkan pada $\ocint{-\pi,\pi}$, dan $\sequencen{s_n}$ merupakan barisan
dari jumlah Fourier.    Misalkan $x\in\ooint{-\pi,\pi}$ adalah
sedemikian sehingga

(i) ada $a\in\Bbb C$ sedemikian rupa

\inset{{\it atau}
$\int_{-\pi}^x\Bover{a-f(t)}{x-t}dt$ ada di $\Bbb C$}

\inset{{\it atau} ada beberapa $\delta>0$ sehingga $f$ merupakan variasi
yang dibatasi pada $[x-\delta,x]$, dan $a=\lim_{t\in\dom f,t\uparrow x}f(t)$}

(ii) ada $b\in\Bbb C$ sedemikian rupa

\inset{{\it atau}
$\int_x^{\pi}\Bover{f(t)-b}{t-x}dt$ ada di $\Bbb C$}

\inset{{\it atau} ada beberapa $\delta>0$ sehingga $f$ merupakan variasi
yang dibatasi pada $[x,x+\delta]$, dan $b=\lim_{t\in\dom f,t\downarrow
x}f(t)$.}

\noindent Tunjukkan bahwa $\lim_{n\to\infty}s_n(x)={1\over 2}(a+b)$.
Modifikasi apa yang sesuai untuk mendapatkan batas
$\lim_{n\to\infty}s_n(\pi)$?
%282Xi, 282O

\sqheader 282Xq Misalkan $f$, $g$ adalah fungsi
bernilai kompleks yang dapat diintegralkan pada $\ocint{-\pi,\pi}$, dan
$\pmb{c}=\langle c_k\rangle_{k\in\Bbb Z}$,
$\pmb{d}=\langle d_k\rangle_{k\in\Bbb Z}$
merupakan barisan koefisien Fourier-nya.   Misalkan {\it salah satu} dari
$\sum_{k=-\infty}^{\infty}|c_k|<\infty$ {\it atau}
$\sum_{k=-\infty}^{\infty}(|c_k|^2+|d_k|^2)<\infty$.   Tunjukkan bahwa barisan
koefisien Fourier dari $f\times g$ hanyalah konvolusi
$\pmb{c}*\pmb{d}$ dari $\pmb{c}$ dan $\pmb{d}$ (255Xk).
%282Q

\spheader 282Xr Di 282Ra, apa yang terjadi jika $f(\pi)\ne f(-\pi)$?
%282R

\spheader 282Xs  Misalkan $\langle c_k\rangle_{k\in\Bbb Z}$
adalah barisan bilangan kompleks berujung ganda sehingga
$\sum_{k=-\infty}^{\infty}|kc_k|<\infty$.   Tunjukkan bahwa
$f(x)=\sum_{k=-\infty}^{\infty}c_ke^{ikx}$ ada untuk setiap
$x\in\Bbb R$ dan $f$ terdiferensiasi di mana saja.
%282R

\spheader 282Xt Misalkan $\langle c_k\rangle_{k\in\Bbb Z}$ menjadi barisan bilangan kompleks berujung ganda
sedemikian rupa sehingga
$\sup_{k\in\Bbb Z}|kc_k|<\infty$.   Tunjukkan bahwa terdapat fungsi
yang terintegralkan-kuadrat $f$
pada $\ocint{-\pi,\pi}$ sehingga $c_k$ adalah koefisien Fourier
dari $f$, bahwa $f$ adalah limit hampir di mana-mana dari jumlah Fourier, dan bahwa $f*f*f$ dapat terdiferensiasi.   \Hint{menggunakan
282K untuk menunjukkan bahwa terdapat $f$, dan 282Xn untuk menunjukkan bahwa jumlah Fourier
nya konvergen di setiap titik tempat jumlah Fej\'er-nya konvergen; gunakan 282Q dan 282Xs untuk
menunjukkan bahwa $f*f*f$ terdiferensiasi.}
%282Xs, 282Q

\leader{282Y}{Latihan lebih lanjut (a)}
%\spheader 282Ya
Misalkan $f$ adalah fungsi yang dapat diintegralkan nonnegatif pada
$\ocint{-\pi,\pi}$, dengan koefisien Fourier
$\langle c_k\rangle_{k\in\Bbb Z}$.   Tunjukkan bahwa

\Centerline{$\sum_{j=0}^n\sum_{k=0}^na_j\bar a_kc_{j-k}\ge 0$}

\noindent untuk semua bilangan kompleks $a_0,\ldots,a_n$.   (Lihat juga 285Xu
di bawah.)
%282A

\spheader 282Yb Misalkan $f:\ocint{-\pi,\pi}\to\Bbb C$, $K\ge 0$, $\alpha>0$
menjadi $|f(x)-f(y)|\le K|x-y|^{\alpha}$ untuk semua $x$,
$y\in\ocint{-\pi,\pi}$.   Misalkan $c_k$, $s_n$ menjadi koefisien Fourier
dan jumlah dari $f$.   (i) Tunjukkan bahwa
$\sup_{k\in\Bbb Z}|k|^{\alpha}|c_k|<\infty$.   \Hint{menunjukkan bahwa
$c_k=\bover1{4\pi}\int_{-\pi}^{\pi}(f(x)-\tilde f(x+\bover{\pi}k))
e^{-ikx}dx$.}   (ii) Tunjukkan bahwa jika $f(\pi)=\lim_{x\downarrow-\pi}f(x)$
maka $s_n\to f$ seragam.   (Bandingkan 282Xj.)
%282E

\spheader 282Yc Misalkan $f$ adalah fungsi bernilai kompleks
yang terukur pada $\ocint{-\pi,\pi}$, dan anggaplah
$p\in\coint{1,\infty}$ sedemikian rupa sehingga
$\int_{-\pi}^{\pi}|f|^p<\infty$.   Misalkan
$\sequence{m}{\sigma_m}$ menjadi barisan jumlah Fej\'er dari $f$.   Tunjukkan
bahwa $\lim_{m\to\infty}\int_{-\pi}^{\pi}|f-\sigma_m|^p=0$.
\Hint{menggunakan 245Xl, 255Yk dan ide-ide di 282Ib.}
%282I

\spheader 282Yd Buatlah fungsi kontinu
$h:[-\pi,\pi]\to\Bbb R$ sedemikian rupa sehingga $h(\pi)=h(-\pi)$ tetapi jumlah Fourier
dari $h$ tidak dibatasi di $0$, sebagai berikut.   Tetapkan
$\alpha(m,n)=\int_0^{\pi}
\bover{\sin(m+{1\over2})t\sin(n+{1\over 2})t}{\sin{1\over 2}t}dt$.
Tunjukkan bahwa $\lim_{n\to\infty}\alpha(m,n)=0$
untuk setiap $m$, tetapi $\lim_{n\to\infty}\alpha(n,n)=\infty$.   Tetapkan
$h_0(x)=\sum_{k=0}^{\infty}\delta_k\sin(m_k+{1\over 2})x$ untuk
$0\le x\le\pi$, $0$ untuk $-\pi\le x\le 0$, di mana
$\delta_k>0$, $m_k\in\Bbb N$ sedemikian rupa sehingga
($\alpha$) $\delta_k\le 2^{-k}$, $\delta_k|\alpha(m_k,m_n)|\le 2^{-k}$
untuk setiap $n<k$ (memilih $\delta_k$) ($\beta$)
$\delta_k\alpha(m_k,m_k)\ge k$, $\delta_n|\alpha(m_k,m_n)|\le 2^{-n}$
untuk setiap $n<k$ (memilih
$m_k$).   Sekarang modifikasi $h_0$ di $\coint{-\pi,0}$ dengan menambahkan fungsi variasi terbatas.
%282*

\spheader 282Ye(i) Tunjukkan bahwa $\lim_{n\to\infty}\int_{-\pi}^{\pi}
|\bover{\sin(n+{1\over2})t}{\sin{1\over 2}t}|dt=\infty$.  \Hint{282Xm.}
(ii) Tunjukkan bahwa untuk setiap $\delta>0$ terdapat $n\in\Bbb N$, $f\ge 0$ seperti
sehingga $\int_{-\pi}^{\pi}f\le\delta$, $\int_{-\pi}^{\pi}|s_n|\ge 1$, di mana
$s_n$ adalah jumlah Fourier ke-$n$ dari $f$.   \Hint{ambil $n$ sedemikian rupa sehingga
$\Bover1{2\pi}\int_{-\pi}^{\pi}
|\bover{\sin(n+{1\over2})t}{\sin{1\over 2}t}|dt>\Bover1{\delta}$ dan tetapkan
$f(x)=\Bover{\delta}{\eta}$ untuk $0\le x\le\eta$, $0$ sebaliknya, jika
$\eta$ kecil.}   (iii) Tunjukkan bahwa terdapat fungsi
$f:\ocint{-\pi,\pi}\to\Bbb R$ yang dapat diintegralkan sehingga $\sup_{n\in\Bbb N}\|s_n\|_1$ adalah
tak terbatas, dengan $\sequencen{s_n}$ adalah barisan jumlah Fourier dari
$f$.   \Hint{pengetahuan tentang `Teorema Keterbatasan Seragam' dari analisis fungsional
akan membantu, tetapi $f$ juga dapat dibuat dengan tangan kosong dengan metode
dari 282Yd.}
%282*

\spheader 282Yf Misalkan $u:[-\pi,\pi]\to\Bbb R$ adalah fungsi
yang kontinu mutlak sehingga $u(\pi)=u(-\pi)$ dan $\int_{-\pi}^{\pi}u=0$.   Tunjukkan
bahwa $\|u\|_2\le\|u'\|_2$.   (Ini adalah ketidaksetaraan {\bf Wirtinger}.)
%282J

\spheader 282Yg Untuk $0\le r<1$, $t\in\Bbb R$ mengatur
$A_r(t)=\Bover{1-r^2}{1-2r\cos t+r^2}$.
($A_r$ adalah {\bf kernel Poisson}; lihat 478Xl\Latereditions\ di
Volume 4.)
(i) Tunjukkan $\Bover1{2\pi}\int_{-\pi}^{\pi}A_r=1$.
(ii) Untuk fungsi bernilai real $f$ yang
dapat diintegralkan pada $\ocint{-\pi,\pi}$, dengan koefisien Fourier real
$a_k$, $b_k$ (282Ba), tetapkan
$S_r(x)=\Bover12a_0+\sum_{k=1}^{\infty}r^k(a_k\cos kx+b_k\sin kx)$ untuk
$x\in\ocint{-\pi,\pi}$, $r\in\coint{0,1}$.   Tunjukkan bahwa
$S_r(x)=\Bover1{2\pi}\int_{-\pi}^{\pi}A_r(x-t)f(t)dt$ untuk setiap
$x\in\ocint{-\pi,\pi}$.   \Hint{$A_r(t)=1+2\sum_{n=1}^{\infty}r^n\cos nt$.}
(iii) Tunjukkan bahwa $\lim_{r\uparrow 1}S_r(x)=f(x)$ untuk setiap
$x\in\ooint{-\pi,\pi}$ yang ada di himpunan Lebesgue $f$.
\Hint{223Yg.}  (iv) Tunjukkan bahwa
$\lim_{r\uparrow 1}\int_{-\pi}^{\pi}|S_r-f|=0$.   (v) Tunjukkan bahwa jika $f$
didefinisikan di seluruh $\ocint{-\pi,\pi}$ dan kontinu, dan
$f(\pi)=\lim_{x\downarrow-\pi}f(x)$, maka
$\lim_{r\uparrow 1}\sup_{x\in\ocint{-\pi,\pi}}|S_r(x)-f(x)|=0$.
%282+
}%end of exercises

\endnotes{
\Notesheader{282} Ini adalah bagian yang panjang dengan kumpulan hasil
yang berpotensi membingungkan, jadi mungkin saya harus merekapitulasi.   Terkait dengan fungsi dapat diintegralkan apa pun di
$\ocint{-\pi,\pi}$, kami memiliki jumlah Fourier yang sesuai, yaitu jumlah parsial
yang simetris $\sum_{k=-n}^nc_ke^{ikx}$ dari deret kompleks
$\sum_{k=-\infty}^{\infty}c_ke^{ikx}$, atau, sama halnya, jumlah parsial
${1\over 2}a_0+\sum_{k=1}^na_k\cos kx+b_k\sin kx$ dari deret real
${1\over 2}a_0+\sum_{k=1}^{\infty}a_k\cos kx+b_k\sin kx$.   Koefisien Fourier
 $c_k$, $a_k$, $b_k$ adalah satu-satunya yang alami, karena
jika deret tersebut ingin konvergen dengan keteraturan apa pun maka

\Centerline{$\Bover1{2\pi}\int_{-\pi}^{\pi}
\bigl(\sum_{k=-\infty}^{\infty}c_ke^{ikx}\bigr)e^{-ilx}dx$}

\noindent seharusnya secara bersamaan

\Centerline{$\sum_{k=-\infty}^{\infty}\Bover1{2\pi}\int_{-\pi}^{\pi}
c_ke^{ikx}e^{-ilx}dx=c_l$}

\noindent dan

\Centerline{$\Bover1{2\pi}\int_{-\pi}^{\pi}f(x)e^{-ilx}dx$.}

\noindent (Bandingkan perhitungan di 282J.) Efek dari mengambil jumlah
Fej\'er $\sigma_m(x)$ daripada jumlah Fourier $s_n(x)$ adalah untuk
memuluskan urutannya;   ingat jika $\lim_{n\to\infty}s_n(x)=c$
maka $\lim_{m\to\infty}\sigma_m(x)=c$, oleh 273Ca di bab terakhir.

Sebagian besar pekerjaan di atas berkaitan dengan pertanyaan kapan jumlah Fourier atau
Fej\'er konvergen, dalam arti tertentu, ke fungsi asli $f$.   Seperti yang telah terjadi pada
sebelumnya, di \S245 dan di tempat lain, kami memiliki lebih dari satu jenis konvergensi
yang perlu dipertimbangkan. Konvergensi   {\it norma}, untuk
$\|\,\|_1$ atau $\|\,\|_2$ atau $\|\,\|_{\infty}$, adalah yang paling sederhana;  tiga teorema
 282G, 282Ib dan 282J setidaknya
relatif mudah.   (Saya telah memberikan 282Ib sebagai akibat wajar dari
282Ia; tetapi ada bukti yang lebih mudah dari 282G.   Lihat 282Xd.)
Secara berurutan, kita mempunyai

\inset{jika $f$ kontinu (dan cocok pada $\pm\pi$, yaitu
$f(\pi)=\lim_{t\downarrow-\pi}f(t)$) maka $\sigma_m\to
f$ seragam, yaitu untuk $\|\,\|_{\infty}$ (282G);}

\inset{jika $f$ adalah fungsi apa pun yang dapat diintegralkan, maka $\sigma_m\to f$ untuk
$\|\,\|_1$ (282Ib);}

\inset{jika $f$ adalah fungsi yang terintegralkan-kuadrat, maka $s_n\to f$ untuk
$\|\,\|_2$ (282J);}

\inset{jika $f$ kontinu dan variasi
terbatas (dan cocok di $\pm\pi$), maka $s_n\to f$ seragam (282O).}

\noindent Ada beberapa hasil serupa untuk $\|\,\|_p$ lainnya (282Yc);
tetapi perhatikan bahwa jumlah Fourier tidak perlu konvergen untuk $\|\,\|_1$ (282Ye).

Konvergensi {\it titik demi titik} lebih sulit.   Hasil-hasil yang saya berikan adalah

\inset{jika $f$ adalah fungsi apa pun yang dapat diintegralkan, maka $\sigma_m\to f$ hampir
di mana-mana (282Ia);}

\noindent ini bergantung pada beberapa perhitungan yang cermat di 282H, dan juga pada
hasil mendalam 223D.   Selanjutnya kita memperoleh hasil yang meninjau rata-rata
limit $f$ dari kedua sisi.   Misalkan saya menulis

\Centerline{$f^{\pm}(x)=\Bover12(\lim_{t\uparrow
x}f(t)+\lim_{t\downarrow x}f(t))$}

\noindent kapan pun nilai ini terdefinisi, dengan mengambil $f^{\pm}(\pi)
=\bover12(\lim_{t\uparrow\pi}f(t)+\lim_{t\downarrow-\pi}f(t))$.   Maka
kita mempunyai

\inset{jika $f$ adalah fungsi apa pun yang dapat diintegralkan, $\sigma_m\to f^{\pm}$ di mana pun
$f^{\pm}$ didefinisikan (282I);}

\inset{jika $f$ memiliki variasi terbatas, $s_n\to f^{\pm}$ di mana-mana
(282O).}

\noindent Tentu saja ini berlaku pada titik mana pun di mana $f$ kontinu,
dalam hal ini $f(x)=f^{\pm}(x)$.
Hasil lain dari jenis ini adalah

\inset{jika $f$ adalah fungsi apa pun yang dapat diintegralkan, $s_n\to f$ pada setiap titik tempat
$f$ terdiferensiasi (282L);}

\noindent sebenarnya, hasil ini dapat diperluas secara berguna dengan sedikit sekali pekerjaan tambahan
(282Xi, 282Xp).

Saya tidak dapat mengakhiri daftar ini tanpa menyebutkan teorema yang {\it tidak} saya berikan.   Inilah {\bf teorema Carleson}:

\inset{jika $f$ terintegralkan-kuadrat, $s_n\to f$ hampir di mana-mana}

\noindent ({\smc Carleson 66}).   Saya akan membahasnya di \S286.   Terdapat
kasus khusus elementer di 282Xt.   Hasil ini sebenarnya
valid untuk banyak $f$ lainnya (lihat catatan untuk \S286).

Kekosongan mencolok berikutnya dalam uraian ini adalah tidak adanya contoh
yang menunjukkan sejauh mana hasil-hasil ini merupakan hasil terbaik yang mungkin.   Memang, tidak ada alasan untuk menduga
bahwa terdapat syarat perlu dan cukup
yang alami bagi

\inset{$s_n\to f$ di setiap titik.}

\noindent Namun demikian, kita perlu berupaya menemukan fungsi
kontinu
yang tidak memiliki sifat ini, dan konstruksi contoh oleh du
Bois-Reymond ({\smc Bois-Reymond 1876}) merupakan momen penting dalam sejarah analisis, antara lain karena memaksa para matematikawan
menyadari bahwa beberapa anggapan yang nyaman tentang klasifikasi fungsi
-- pada dasarnya, bahwa fungsi
itu `baik' atau sedemikian buruk sehingga tidak perlu dipedulikan -- ternyata keliru.  Contoh tersebut
bersifat mendidik, tetapi terpaksa saya hilangkan karena keterbatasan ruang; saya memberikan garis besar
salah satu metode yang mungkin di 282Yd.   (Konstruksi terperinci dapat ditemukan
dalam {\smc K\"orner 88}, bab 18, dan bukti keberadaan fungsi
semacam itu dalam {\smc Dudley 89}, 7.4.3.) Jika kita mengizinkan fungsi umum
yang dapat diintegralkan, kita dapat memperoleh hasil yang jauh lebih baik, atau barangkali seharusnya saya
mengatakan jauh lebih buruk; terdapat $f$ yang dapat diintegralkan sedemikian sehingga $\sup_{n\in\Bbb
N}|s_n(x)|=\infty$ untuk setiap $x\in\ocint{-\pi,\pi}$
({\smc Kolmogorov 1926}; lihat {\smc Zygmund 59}, $\S\S$VIII.3-4).

Di 282C saya menyebutkan dua jenis masalah. Yang pertama -- kapan deret
Fourier
dapat dijumlahkan? -- setidaknya telah dibahas panjang lebar, meskipun saya
tidak dapat mengaku telah memberikan lebih dari sekadar contoh dari apa yang diketahui.
Yang kedua -- bagaimana sifat-sifat $c_k$ mencerminkan sifat-sifat $f$? -- nyaris belum saya sentuh.   Saya memberikan apa yang menurut saya merupakan tiga hasil
yang paling penting di bidang ini.   Yang pertama adalah

\inset{jika $f$ dan $g$ mempunyai koefisien Fourier yang sama, maka keduanya sama
hampir di mana-mana (282Ic).}

\noindent Ini setidaknya memberi tahu kita bahwa pada prinsipnya kita harus dapat
mempelajari hampir semua hal tentang $f$ dengan melihat deret Fourier-nya.
(Misalnya, 282Ya menjelaskan kondisi perlu dan cukup agar
$f$ menjadi nonnegatif hampir di mana-mana.) Yang kedua adalah

\inset{$f$ terintegralkan-kuadrat jika dan hanya jika
$\sum_{k=-\infty}^{\infty}|c_k|^2<\infty$;}

\noindent sebenarnya,

\inset{$\sum_{k=-\infty}^{\infty}|c_k|^2
=\Bover1{2\pi}\int_{-\pi}^{\pi}|f|^2$
(282J).}

\noindent Tentu saja ini bersifat mendasar, karena menunjukkan bahwa koefisien
Fourier memberikan isomorfisme ruang Hilbert alami antara $L^2$
dan $\ell^2$ (282K). Perlu saya tambahkan bahwa meskipun ruang Hilbert real
$L^2_{\Bbb R}$ dan $\ell^2_{\Bbb R}$ isomorfik sebagai
ruang hasil kali dalam (282Xg), keduanya jelas tidak isomorfik sebagai kisi Banach;
misalnya, $\ell^2_{\Bbb R}$ memiliki elemen `atomik'
$\pmb{c}$ sedemikian rupa sehingga jika $0\le\pmb{d}\le\pmb{c}$ maka $\pmb{d}$ adalah
kelipatan $\pmb{c}$, sedangkan $L^2_{\Bbb R}$ tidak memilikinya.   Barangkali yang
lebih penting lagi adalah

\inset{koefisien Fourier dari konvolusi $f*g$ hanyalah suatu kelipatan
skalar dari
hasil kali koefisien Fourier $f$ dan $g$ (282Q);}

\noindent tetapi untuk menggunakannya secara efektif kita perlu mempelajari struktur
aljabar Banach $L^1$, dan saya tidak punya pilihan selain segera meninggalkan
jalur ini. (Pembahasannya akan menjadi bagian penting Bab 44 dalam
Volume 4.)   282Xt memberikan suatu
konsekuensi elementer, sedangkan 282Xq memberikan uraian yang sangat parsial mengenai
hubungan antara hasil kali $f\times g$ dari dua fungsi dan
hasil kali konvolusi barisan koefisien Fourier keduanya.

Jumlah Fej\'er yang dibahas dalam bagian ini merupakan salah satu cara untuk mengatasi
kesulitan konvergensi yang berkaitan dengan jumlah Fourier.   Ketika
kita membahas transformasi Fourier dalam dua bagian berikutnya, kita akan memerlukan
beberapa siasat tambahan.   Jenis penghalusan yang berbeda
diperoleh dengan menggunakan kernel Poisson sebagai pengganti kernel Dirichlet atau Fej\'er
(282Yg).

Saya mengakhiri catatan ini dengan komentar mengenai bilangan $2\pi$. Bilangan ini muncul dalam
hampir setiap rumus yang melibatkan deret Fourier, tetapi menurut saya dapat
dihilangkan sepenuhnya dari bagian ini, setidaknya, dengan menormalisasi ulang
ukuran pada $\ocint{-\pi,\pi}$. Jika alih-alih ukuran Lebesgue
$\mu$ kita menggunakan ukuran $\nu=\bover1{2\pi}\mu$ di seluruh pembahasan, setiap faktor
$2\pi$ akan hilang. (Bandingkan keterangan di 282Bb mengenai
kemungkinan menghitung integral pada $S^1$.) Namun menurut saya kebanyakan dari kita
lebih suka mengingat letak faktor $2\pi$ dalam setiap rumus daripada
berurusan dengan ukuran yang tidak lazim.
}%end of notes

\discrpage
